\begin{exname}Problem\end{exname}
% begin pdf problem style
\renewcommand{\ocExName}{Problem}
% end pdf problem style
\chapter{Almost every number}\label{ch:harder}

\begin{center}\itshape
This conjecture, posed in 1941, can be thought of as the ultimate generalization of Khintchine's Theorem, describing how well a
typical real number can be approximated by rational ones.\\
\upshape --- International Mathematical Union, citation for James Maynard's Fields Medal (2022)
\end{center}

In 1924 Aleksandr Khinchin settled a question that sounds as if it had no answer: how well can a \emph{typical} real number be
approximated by fractions? Dirichlet had shown that every irrational $x$ has infinitely many fractions with $|x-\frac pq|<\frac1{q^2}$,
and Liouville that some numbers, his own constant among them, are approximated absurdly well. Khinchin replaced ``every'' and
``some'' by ``almost every'', in the sense of Lebesgue measure, and found a clean law. If $q^2\psi(q)$ decreases, then almost every
$x$ has infinitely many fractions with $|x-\frac pq|<\psi(q)$ when $\sum q\psi(q)$ diverges, and almost none does when it converges. So
almost every number is approximated to within $\frac1{q^2\log q}$ infinitely often, and almost none to within
$\frac1{q^2(\log q)^{1.01}}$. The boundary between the two sides is as sharp as the boundary between a divergent and a convergent series.

Khinchin's monotonicity assumption looked like a technical convenience. In 1941 Richard Duffin and Albert Schaeffer showed that it is
not: without it the theorem is false, because a fraction $\frac pq$ that is not in lowest terms is counted again under a smaller
denominator. They proposed the right statement, counting only reduced fractions and replacing $q\psi(q)$ by $\varphi(q)\psi(q)$, with
Euler's function. In 1961 Patrick Gallagher showed that the set of good $x$ always has measure $0$ or $1$, so only the
alternative was left to decide. It took another fifty-eight years. In July 2019 Dimitris Koukoulopoulos and James Maynard posted a
proof of the Duffin--Schaeffer conjecture; it appeared in the \emph{Annals of Mathematics} in 2020, and it is one of the results the
International Mathematical Union cited when it gave Maynard the Fields Medal in 2022.

This chapter is a set of problems about the machinery behind such statements, and all of it is in this volume: the Borel--Cantelli
lemma of Chapter~\ref{ch:measures}, Baire's theorem of Chapter~\ref{ch:baire}, Lebesgue's density theorem of Chapter~\ref{ch:radon} and the
Hausdorff dimension of Chapter~\ref{ch:metricmeasures}. Continued fractions come from Volume~2, Chapter~9. Nothing is borrowed from a later
volume. The classical proofs of several results below call on the ergodic theorem, which is proved in Volume~11; here a cruder estimate,
proved on the spot, does the same work, and the sharp constants wait for that volume. The first section
proves Khinchin's theorem. The second measures the exceptional sets that Lebesgue measure cannot see, with Hausdorff dimension. The
third shows that almost every $x$ makes almost any sequence $(a_nx)$ equidistributed, and the last compares measure with category,
which disagree about what ``typical'' means.

\begin{think}
Pick $x$ at random in $[0,1]$. Do you expect $|x-\frac pq|<\frac1{q^2\log q}$ to have infinitely many solutions? Now take $x=e$, whose
continued fraction $[2;1,2,1,1,4,1,1,6,\dots]$ has small, regular digits. Is $e$ better or worse approximated than a random number?
\end{think}

Throughout, $x=[a_0;a_1,a_2,\dots]$ is the continued fraction of an irrational $x$, with complete quotients $x_{n+1}=[a_{n+1};a_{n+2},\dots]$
and convergents $\frac{p_n}{q_n}$, as in Volume~2, Chapter~9. For $x\in(0,1)$ the \emph{cylinder} $I(a_1,\dots,a_n)$ is the set of irrational $x$ whose
first $n$ digits are $a_1,\dots,a_n$. Its closure is the interval with ends $\frac{p_n}{q_n}$ and $\frac{p_n+p_{n-1}}{q_n+q_{n-1}}$, of length $\frac1{q_n(q_n+q_{n-1})}$;
Problem~\ref{pb:borelbernstein}(a) proves this. The map $G(x)=\{\frac1x\}$ is the
\emph{Gauss map}: it deletes the first digit, so that $a_{n+1}(x)=a_1(G^nx)$. $\lambda$ is Lebesgue measure, $\|t\|$ is the distance from $t$ to the nearest integer, $e(t)=e^{2\pi it}$, and $\dim_H$ is
Hausdorff dimension (Chapter~\ref{ch:metricmeasures}). ``Borel--Cantelli'' always means Proposition~\ref{prop:measureprops}(e). We use freely the following facts of Volume~2, Theorem~9.2 and Lemma~9.4:
\[
\Bigl|x-\frac{p_n}{q_n}\Bigr|=\frac1{q_n(x_{n+1}q_n+q_{n-1})},\qquad\text{and}\qquad |qx-p|\ge|q_nx-p_n|\ \text{ whenever }1\le q<q_{n+1}.
\]
In particular $\frac1{(a_{n+1}+2)q_n^2}<|x-\frac{p_n}{q_n}|<\frac1{a_{n+1}q_n^2}$.

\section{Approximation, almost surely}

The convergence half of Khinchin's theorem is one line of Borel--Cantelli; the divergence half is where the work is, because the events
``$x$ is close to some $\frac pq$'' for different $q$ are not independent. Continued fractions repair this: the digits of a random number
are nearly independent (Problem~\ref{pb:borelbernstein}(a)), and the size of the digit $a_{n+1}$ is exactly the quality of the approximation
$\frac{p_n}{q_n}$.

\begin{exercise}[$\star$]\label{pb:khinconv}
Let $\psi:\N^*\to[0,\infty[$ with $\sum_qq\psi(q)<\infty$.
\begin{enumerate}[label=(\alph*)]
\item Show that for almost every $x\in\R$ the inequality $|x-\frac pq|<\psi(q)$ has only finitely many solutions $(p,q)\in\Z\times\N^*$.
\item Deduce that for every $\varepsilon>0$, almost every $x$ satisfies $|x-\frac pq|<q^{-2-\varepsilon}$ only finitely often, and recover the
fact that the Liouville numbers form a null set.
\item Show that for almost every $x$ there is $c(x)>0$ with $|x-\frac pq|\ge\frac{c(x)}{q^2\log^2q}$ for all $p$ and all $q\ge2$.
\end{enumerate}
\end{exercise}

\begin{exercise}[$\star$]\label{pb:badnull}
Call $x$ \emph{badly approximable} if it is irrational and there is $c>0$ with $|x-\frac pq|\ge\frac c{q^2}$ for all $p$ and $q\ge1$; let
$\mathrm{Bad}$ be the set of such numbers.
\begin{enumerate}[label=(\alph*)]
\item Show that $\mathrm{Bad}$ is meagre.
\item Show that $x$ is badly approximable if and only if its digits $a_n$ are bounded.
\item Deduce that $\mathrm{Bad}$ is null.
\item Show that $\mathrm{Bad}$ is nevertheless dense and uncountable.
\end{enumerate}
\end{exercise}

\begin{exercise}[$\star\star$]\label{pb:borelbernstein}
\emph{The Borel--Bernstein theorem.} Let $\varphi:\N^*\to]0,\infty[$.
\begin{enumerate}[label=(\alph*)]
\item Let $I=I(a_1,\dots,a_n)$ be a cylinder and $B\subseteq(0,1)$ a Borel set. Find the endpoints of $I$ and show that its length is $\frac1{q_n(q_n+q_{n-1})}$. Show that
\[
\frac12\lambda(B)\le\frac{\lambda(I\cap G^{-n}B)}{\lambda(I)}\le2\lambda(B).
\]
In particular $\frac1{2k}\le\lambda(a_{n+1}\ge k\mid I)\le\frac2k$ for every integer $k\ge1$.
\item Show that if $\sum_n\frac1{\varphi(n)}<\infty$, then for almost every $x\in(0,1)$, $a_n(x)\ge\varphi(n)$ for only finitely many $n$.
\item Show that if $\sum_n\frac1{\varphi(n)}=\infty$, then for almost every $x$, $a_n(x)\ge\varphi(n)$ for infinitely many $n$.
\item Deduce that for almost every $x$, $\limsup_n\frac{a_n(x)}{n\log n}=\infty$, while $a_n(x)<n(\log n)^{1.01}$ for all large $n$.
\end{enumerate}
\end{exercise}

\begin{exercise}[$\star\star\star$]\label{pb:khinchin}
\emph{Khinchin's theorem.} Let $\psi:\N^*\to[0,\infty[$ be such that $f(q)=q^2\psi(q)$ is non-increasing. Let $W(\psi)$ be the set of
$x\in(0,1)$ for which $|x-\frac pq|<\psi(q)$ has infinitely many solutions.
\begin{enumerate}[label=(\alph*)]
\item Show that $\lambda(W(\psi))=0$ if $\sum_qq\psi(q)<\infty$.
\item Show that $q_n\le\prod_{i\le n}(a_i+1)$, and that $\int_0^1\prod_{i\le n}(a_i(x)+1)^{1/2}\dd x\le c^n$ with
$c=2\sum_{k\ge1}\frac1{k\sqrt{k+1}}<4.4$. Deduce that for almost every $x$, $q_n(x)\le e^{5n}$ for all large $n$.
\item Let $x$ be irrational and $n\ge1$. Show that if $a_{n+1}(x)\ge\frac1{f(q_n)}$, then $|x-\frac{p_n}{q_n}|<\psi(q_n)$.
\item Let $Q_n=\lfloor e^{5n}\rfloor$. Show that $\sum_qq\psi(q)=\infty$ implies $\sum_nf(Q_n)=\infty$.
\item Deduce that $\lambda(W(\psi))=1$ if $\sum_qq\psi(q)=\infty$.
\item Conclude that almost every $x$ satisfies $|x-\frac pq|<\frac1{q^2\log q}$ for infinitely many fractions, but
$|x-\frac pq|<\frac1{q^2(\log q)^{1+\varepsilon}}$ for only finitely many, for every $\varepsilon>0$.
\end{enumerate}
\end{exercise}

\begin{exercise}[$\star\star$]\label{pb:econt}
\emph{The number $e$ is not typical.} We quote Euler's continued fraction $e=[2;1,2,1,1,4,1,1,6,\dots]$: $a_0=2$, $a_{3k-1}=2k$ for $k\ge1$, and
all other digits equal $1$. (Quoted, not proved in this series so far; a two-page proof is in Cohn's note in the sources. Everything else in this
problem is proved.)
\begin{enumerate}[label=(\alph*)]
\item Show that $2^kk!\le q_{3k}\le 3^{3k}(k+1)!$ for all $k\ge1$, and deduce that $\log q_n\sim\frac n3\log n$.
\item Show that there is $c>0$ with $\bigl|e-\frac pq\bigr|\ge\frac{c\log\log q}{q^2\log q}$ for all $p$ and $q\ge3$. In particular, for every
$\varepsilon>0$, $|e-\frac pq|<q^{-2-\varepsilon}$ has only finitely many solutions.
\item Deduce that $|e-\frac pq|<\frac1{q^2\log q}$ has only finitely many solutions: $e$ lies in the null set of exceptions of
Problem~\ref{pb:khinchin}(f). (Numerically, one finds nine solutions in integers $q\ge2$, the last one $\frac{1084483}{398959}$; eight of them are in lowest terms.)
\end{enumerate}
\end{exercise}

\section{The size of the exceptions}

Lebesgue measure sees the typical number and is blind to the exceptions: the Liouville numbers, the badly approximable numbers and the
numbers with $|x-\frac pq|<q^{-3}$ infinitely often are all null. Hausdorff dimension tells them apart. The tools are a covering, for upper
bounds, and a cleverly spread measure with the mass distribution principle (Lemma~\ref{lem:massdist}), for lower bounds.

\begin{exercise}[$\star$]\label{pb:wtau}
For $\tau>2$ let $W(\tau)$ be the set of $x\in[0,1]$ such that $|x-\frac pq|<q^{-\tau}$ for infinitely many fractions $\frac pq$.
\begin{enumerate}[label=(\alph*)]
\item Show that $\dim_HW(\tau)\le\frac2\tau$.
\item Deduce that the Liouville numbers form a set of Hausdorff dimension $0$.
\end{enumerate}
\end{exercise}

\begin{exercise}[$\star\star\star$]\label{pb:jarnik}
\emph{The Jarník--Besicovitch theorem: $\dim_HW(\tau)=\frac2\tau$.}
\begin{enumerate}[label=(\alph*)]
\item \emph{A Cantor-type lemma.} Let $[0,1]=E_0\supseteq E_1\supseteq\cdots$, where each $E_k$ is a finite union of disjoint closed intervals, each
interval of $E_{k-1}$ contains at least $m_k\ge2$ intervals of $E_k$, and the intervals of $E_k$ are separated by gaps of length at least
$\varepsilon_k$, with $\varepsilon_1>\varepsilon_2>\cdots$ and $m_k\varepsilon_k<1$. Show that $F=\bigcap_kE_k$ satisfies
\[
\dim_HF\ge\liminf_{k\to\infty}\frac{\log(m_1m_2\cdots m_{k-1})}{-\log(m_k\varepsilon_k)}.
\]
\item Let $J\subseteq\,]0,1[$ be an interval of length $L$, and $Q$ large with $LQ\ge2$. Show that $J$ contains at least $\frac{LQ^2}{4\log Q}$
fractions $\frac pq$ with $q$ a prime in $]Q,2Q]$, and that any two of them are at distance at least $\frac1{4Q^2}$.
\item Let $\tau>2$ and $\log Q_k=2^{k^2}$. Build sets $E_k$ as in (a) out of intervals $[\frac pq-\frac12q^{-\tau},\frac pq+\frac12q^{-\tau}]$ with $q$ a
prime in $]Q_k,2Q_k]$, such that $F\subseteq W(\tau)$.
\item Show that the right side of (a) equals $\frac2\tau$ for your construction, and conclude.
\end{enumerate}
\end{exercise}

\begin{exercise}[$\star\star$]\label{pb:good}
Let $D$ be the set of irrational $x\in(0,1)$ whose digits tend to infinity: $a_n(x)\to\infty$.
\begin{enumerate}[label=(\alph*)]
\item Show that $D$ is null, and that it is uncountable.
\item Show that $\dim_HD\le\frac12$. (In fact $\dim_HD=\frac12$, a theorem of Good.)
\end{enumerate}
\end{exercise}

\begin{exercise}[$\star\star\star$]\label{pb:etwo}
Let $E_{\{1,2\}}$ be the set of irrational $x\in(0,1)$ all of whose digits are $1$ or $2$.
\begin{enumerate}[label=(\alph*)]
\item Show that $E_{\{1,2\}}$ is a compact set contained in $\bigl[\frac{\sqrt3-1}2,\sqrt3-1\bigr]$, null and uncountable.
\item Show that $\dim_HE_{\{1,2\}}\le\frac{\log2}{2\log\varphi}=0.7202\ldots$, where $\varphi=\frac{1+\sqrt5}2$.
\item Show that $\dim_HE_{\{1,2\}}\ge\frac{\log2}{2\log(1+\sqrt2)}=0.3932\ldots$
\item Direct computation gives $\dim_HE_{\{1,2\}}=0.5312805062772051\ldots$ as the value of $s$ for which the operator
$\mathcal L_sf(x)=\sum_{a\in\{1,2\}}(a+x)^{-2s}f\bigl(\frac1{a+x}\bigr)$ has spectral radius $1$. Explain why the sum
$\sum_{a_1,\dots,a_n\in\{1,2\}}|I(a_1,\dots,a_n)|^s$ should be compared with the $n$-th power of that spectral radius.
\end{enumerate}
\end{exercise}

\begin{figure}[ht]
\centering
\fig{ch01-etwo}
\caption{The first five generations of the construction of $E_{\{1,2\}}$ inside $[0.36,0.74]$: at generation $n$, the $2^n$ intervals spanned by
$E_{\{1,2\}}\cap I(a_1,\dots,a_n)$, coloured charcoal when $a_n=1$ and rust when $a_n=2$. Each interval splits into two separated by a gap, which is
what Problem~\ref{pb:etwo}(c) needs; the intervals shrink at rates between $\varphi^{-2n}$ and $(1+\sqrt2)^{-2n}$.}\label{fig:etwo}
\end{figure}

\begin{exercise}[$\star\star\star$]\label{pb:eggleston}
\emph{The Besicovitch--Eggleston theorem.} For $0<p<1$ let $B_p$ be the set of $x\in[0,1]$ such that the proportion of ones among the first $n$
binary digits of $x$ tends to $p$, and $H(p)=-p\log p-(1-p)\log(1-p)$.
\begin{enumerate}[label=(\alph*)]
\item Let $\mu$ be a probability on $[0,1]$ and $A$ a Borel set with $\mu(A)>0$, such that for every $x\in A$,
$\liminf_n\frac{-\log_2\mu(D_n(x))}n\ge s$, where $D_n(x)$ is the dyadic interval of length $2^{-n}$ containing $x$. Show that $\dim_HA\ge s$.
\item Show that $\dim_HB_p\ge\frac{H(p)}{\log2}$.
\item Show that $\binom nk\le e^{nH(k/n)}$, and deduce that $\dim_HB_p\le\frac{H(p)}{\log2}$.
\item Compute $\dim_HB_{1/3}$, and explain why $B_{1/3}$ is null although its dimension is close to $1$.
\end{enumerate}
\end{exercise}

\begin{exercise}[$\star$]\label{pb:cantorL}
Let $C$ be the Cantor set.
\begin{enumerate}[label=(\alph*)]
\item Show that $x=\sum_{k\ge1}2\varepsilon_k3^{-k!}$ is a Liouville number for every sequence $(\varepsilon_k)$ in $\{0,1\}$ with infinitely many ones.
\item Deduce that $C$ contains uncountably many Liouville numbers, although $C$ and the Liouville numbers are both null, and that
$C\cap W(\tau)\ne\varnothing$ for every $\tau>2$.
\item What is $\dim_H(C\cap\mathcal L)$, where $\mathcal L$ is the set of Liouville numbers?
\end{enumerate}
\end{exercise}

\section{Equidistribution for almost every \texorpdfstring{$x$}{x}}

Weyl's criterion (Volume~5, Theorem~6.13) turns equidistribution into the decay of exponential sums, and Hilbert space turns the decay of exponential
sums \emph{on average over $x$} into a one-line computation. In his memoir of 1916 Weyl drew the consequence that for \emph{every} sequence of
distinct integers $a_n$, almost every $x$ makes $(a_nx)$ equidistributed. The price of ``almost every'' is that one never knows which $x$.

\begin{exercise}[$\star\star$]\label{pb:weylae}
\emph{Weyl's metric theorem.} Let $(a_n)$ be a sequence of distinct integers. Show that for almost every real $x$, the sequence $(a_nx)$ is equidistributed
modulo $1$. (Show first that $S_N(x)=\frac1N\sum_{n\le N}e(ha_nx)$ satisfies $\int_0^1|S_N|^2=\frac1N$ for $h\ne0$.)
\end{exercise}

\begin{exercise}[$\star$]\label{pb:factorial}
\begin{enumerate}[label=(\alph*)]
\item Show that $(n!\,x)$ is equidistributed modulo $1$ for almost every $x$.
\item Show that $\{n!\,e\}\to0$, where $\{t\}$ is the fractional part.
\item Show that the set of $x$ with $\{n!\,x\}\to0$ is dense, uncountable and null.
\end{enumerate}
\end{exercise}

\begin{exercise}[$\star\star\star$]\label{pb:del}
\emph{The Davenport--Erdős--LeVeque criterion.} Let $s_n:[a,b]\to\R$ be measurable functions, $S_N(x)=\frac1N\sum_{n\le N}e(s_n(x))$ and
$I(N)=\int_a^b|S_N(x)|^2\dd x$. Suppose that $\sum_{N\ge1}\frac{I(N)}N<\infty$.
\begin{enumerate}[label=(\alph*)]
\item Let $\rho>1$. Show that one can choose integers $N_1<N_2<\cdots$ with $\rho^m\le N_m<\rho^{m+1}$ for all large $m$ and $\sum_mI(N_m)<\infty$.
\item Show that $S_N(x)\to0$ for almost every $x\in[a,b]$.
\item Compare with the argument of Problem~\ref{pb:weylae}: what does (b) gain?
\end{enumerate}
\end{exercise}

\begin{exercise}[$\star\star\star$]\label{pb:koksma}
\emph{Koksma's theorem.} Show that for almost every $x>1$, the sequence $(x^n)$ is equidistributed modulo $1$. (For $m\ne n$ and $h\ne0$, bound
$\int_a^be\bigl(h(x^m-x^n)\bigr)\dd x$ on an interval $[a,b]\subseteq\,]1,\infty[$.) Name some $x>1$ for which $(x^n)$ is \emph{not} equidistributed.
\end{exercise}

\begin{figure}[ht]
\centering
\fig{ch01-koksma}
\caption{Histograms of $\{x^n\}$ for $1\le n\le 4000$ and two values of $x$, in $20$ bins of width $0.05$ (expected count $200$ per bin): $x=1.37$ (charcoal),
whose powers look equidistributed as Koksma's theorem predicts for almost every $x$, and the golden ratio $x=\varphi$ (rust), for which $\{\varphi^n\}$
accumulates at $0$ and $1$ because $\varphi^n+(-\varphi)^{-n}$ is an integer. The counts for $\varphi$ in the two end bins, $1997$ each, are cut off.}\label{fig:koksma}
\end{figure}

\begin{exercise}[$\star\star$]\label{pb:cantorfourier}
\emph{The Cantor measure.} Let $\mu$ be the law of $X=\sum_{k\ge1}2\varepsilon_k3^{-k}$, where the $\varepsilon_k$ are independent with
$P(\varepsilon_k=0)=P(\varepsilon_k=1)=\frac12$, and $\hat\mu(t)=\int e(-tx)\dd\mu(x)$.
\begin{enumerate}[label=(\alph*)]
\item Show that $\hat\mu(t)=e\bigl(-\frac t2\bigr)\prod_{k\ge1}\cos\frac{2\pi t}{3^k}$.
\item Show that $|\hat\mu(3^m)|=|\hat\mu(1)|\ne0$ for all $m\ge0$, so that $\hat\mu(n)\not\to0$. Why does this not contradict the Riemann--Lebesgue
lemma? (Numerically, one finds $|\hat\mu(1)|=0.37143\ldots$)
\item Show that for \emph{every} $x$ in the Cantor set, the sequence $(3^nx)$ is not equidistributed modulo $1$.
\end{enumerate}
\end{exercise}

\section{Measure against category}

``Typical'' has two meanings on the line: outside a null set, or outside a meagre set (Chapter~\ref{ch:baire}). They are both perfectly reasonable, and
they disagree completely, since $\R$ is the union of a null set and a meagre set (Theorem~\ref{thm:liouville}). The problems of this section play the two
notions against each other.

\begin{exercise}[$\star$]\label{pb:erdosL}
\emph{Erdős's theorem.} Show that every real number is the sum of two Liouville numbers, in uncountably many ways. Compare with
Problem~\ref{pb:wtau}: a set of dimension $0$ whose sumset is the whole line.
\end{exercise}

\begin{exercise}[$\star$]\label{pb:halves}
Let $0<\alpha<1$. Show that there is no Lebesgue measurable set $A\subseteq\R$ with $\lambda(A\cap I)=\alpha\lambda(I)$ for every interval $I$, nor even one
with $\varepsilon\lambda(I)\le\lambda(A\cap I)\le(1-\varepsilon)\lambda(I)$ for every interval $I$, for some fixed $\varepsilon>0$. (Compare with Exercise~\ref{ex:everywhere}.)
\end{exercise}

\begin{exercise}[$\star\star$]\label{pb:similar}
\begin{enumerate}[label=(\alph*)]
\item Let $A\subseteq\R$ be measurable with $\lambda(A)>0$, and $S\subseteq\R$ finite. Show that $A$ contains an \emph{affine copy} $x+tS$ of $S$, with
$t>0$.
\item Let $(r_n)$ be an enumeration of $\Q\cap[0,1]$ and $A=[0,1]\setminus\bigcup_n\,]r_n-4^{-n},r_n+4^{-n}[$. Show that $A$ has positive measure and is
nowhere dense, and deduce that a nowhere dense set may contain affine copies of every finite set.
\item Show that the set of Liouville numbers, which is null, contains an affine copy $x+tS$ of every \emph{countable} set $S$, for every $t>0$.
\end{enumerate}
\end{exercise}

\begin{exercise}[$\star\star$]\label{pb:typicalcf}
Let $X$ be the set of irrational $x\in(0,1)$ with
\[
\liminf_{n\to\infty}(a_1a_2\cdots a_n)^{1/n}=1\qquad\text{and}\qquad\limsup_{n\to\infty}(a_1a_2\cdots a_n)^{1/n}=\infty.
\]
\begin{enumerate}[label=(\alph*)]
\item Show that $X$ is comeagre in $[0,1]$.
\item Show that $X$ is null: the typical number of Baire is not the typical number of Lebesgue.
\end{enumerate}
\end{exercise}

\begin{exercise}[$\star\star\star$]\label{pb:bernsteinset}
\emph{A Bernstein set.} Assume the axiom of choice. Write $\mathfrak c=|\R|$.
\begin{enumerate}[label=(\alph*)]
\item Show that every uncountable closed set $F\subseteq\R$ has cardinality $\mathfrak c$. (Remove from $F$ the points having a countable neighbourhood in $F$,
and build a Cantor set inside what remains.)
\item Show that there are exactly $\mathfrak c$ closed subsets of $\R$.
\item Construct a set $A\subseteq\R$ such that every uncountable closed set meets both $A$ and $\R\setminus A$.
\item Show that $A$ is not Lebesgue measurable: every measurable subset of $A$, or of $\R\setminus A$, is null, and the outer measure of $A\cap[0,1]$ is $1$.
\item Show that $A$ is neither meagre nor comeagre in any interval.
\end{enumerate}
\end{exercise}

\section{What we proved, and what comes next}

Every problem of this chapter ends in a statement about almost every number, or about the size of the numbers left out. Khinchin's theorem
(Problem~\ref{pb:khinchin}) says how well almost every number is approximated by fractions, and its proof rests on one estimate, the bounded distortion of
Problem~\ref{pb:borelbernstein}(a), used with Borel--Cantelli in both directions. The Jarník--Besicovitch theorem (Problem~\ref{pb:jarnik}) and the
Besicovitch--Eggleston theorem (Problem~\ref{pb:eggleston}) give the Hausdorff dimension of two families of null sets, by the mass distribution
principle of Chapter~\ref{ch:metricmeasures}. Weyl's metric theorem and Koksma's theorem (Problems~\ref{pb:weylae} and~\ref{pb:koksma}) make
almost every $x$ produce equidistributed sequences. The last section shows that ``almost every'' in the sense of measure and in the sense of
category are different notions, and ends with a set that neither can handle.

Three things were quoted and not proved: Euler's continued fraction of $e$, used in Problem~\ref{pb:econt}; Lévy's constant; and Khinchin's
constant. The last two are named in the leads and used nowhere. They are limits of averages along the orbit of the Gauss map, and the tool that
computes such limits is the ergodic theorem. Volume~11 proves it, and comes back to continued fractions with it. Two natural companions of these problems need it too, and are in that volume: Minkowski's question-mark function, which is continuous
and strictly increasing with derivative $0$ almost everywhere, and the law of the orbit of a typical point of the Cantor set under multiplication by $3$.

The next volume, \emph{Analysis}, takes the spaces $L^p$ of Chapter~\ref{ch:lp} and the measures of this volume and does Fourier analysis with them.

\section*{Hints}

\begin{hint}{pb:khinconv}
Work in $[0,1]$. For each $q$, the $x$ that are within $\psi(q)$ of some $\frac pq$ form at most $q+2$ intervals; apply the first Borel--Cantelli lemma.
\end{hint}

\begin{hint}{pb:badnull}
For (a), the set of $x$ with $|x-\frac pq|\ge\frac c{q^2}$ for all $\frac pq$ is closed and contains no rational. For (b), use the two estimates recalled at the
start of the chapter. For (c), Problem~\ref{pb:borelbernstein}(c) with $\varphi(n)=n$; its proof does not use this problem.
\end{hint}

\begin{hint}{pb:borelbernstein}
For (a), the points of $I$ are $\frac{p_n+tp_{n-1}}{q_n+tq_{n-1}}$ with $t=G^nx$; compute the derivative in $t$. For (c), bound the measure of
$\{a_n<k_n\text{ for }m\le n\le N\}$ by conditioning on the first $N-1$ digits, one step at a time.
\end{hint}

\begin{hint}{pb:khinchin}
For (b), integrate over one cylinder of rank $n-1$ at a time, with the upper bound of Problem~\ref{pb:borelbernstein}(a) for the law of $a_n$
there; then Markov's inequality and Borel--Cantelli. For (d), cut the series $\sum f(q)/q$ into blocks $Q_{n-1}<q\le Q_n$. For (e), apply Problem~\ref{pb:borelbernstein} with
$\varphi(n+1)=1/f(Q_n)$.
\end{hint}

\begin{hint}{pb:econt}
For (a), $a_nq_{n-1}\le q_n\le(a_n+1)q_{n-1}$. For (b), a fraction $\frac pq$ with $q_n\le q<q_{n+1}$ is no better than $\frac{p_n}{q_n}$, and the digit $a_{n+1}$ is
at most of order $n$.
\end{hint}

\begin{hint}{pb:wtau}
Cover $W(\tau)$, for each $Q$, by the intervals of radius $q^{-\tau}$ around the fractions $\frac pq$ with $q\ge Q$.
\end{hint}

\begin{hint}{pb:jarnik}
For (a), give each interval of $E_k$ the mass $(m_1\cdots m_k)^{-1}$, and estimate the mass of a set $U$ with $\varepsilon_k\le|U|<\varepsilon_{k-1}$ in two
ways: it meets at most one interval of $E_{k-1}$, and at most $2|U|/\varepsilon_k$ intervals of $E_k$. For (b), the prime number theorem, and
$|\frac pq-\frac{p'}{q'}|\ge\frac1{qq'}$. For (c), keep only the fractions in the middle half of each interval.
\end{hint}

\begin{hint}{pb:good}
For (a), by Problem~\ref{pb:borelbernstein}(a) the digit $a_n$ equals $1$ with conditional probability at least $\frac14$, whatever the earlier digits. For (b), cover the set of $x$ with $a_n\ge N$ for all $n>m$ by cylinders, using $|I(a_1,\dots,a_n)|\le(a_1\cdots a_n)^{-2}$, and choose $N$ with
$\sum_{a\ge N}a^{-2s}<1$.
\end{hint}

\begin{hint}{pb:etwo}
For (a), Problem~\ref{pb:badnull}. For (b), $q_n$ is at least a Fibonacci number. For (c), apply Problem~\ref{pb:jarnik}(a) to the intervals spanned by $E_{\{1,2\}}\cap I(a_1,\dots,a_n)$:
the gap between the two children is the image of a fixed gap under the map $t\mapsto\frac{p_n+tp_{n-1}}{q_n+tq_{n-1}}$, and $q_n\le(1+\sqrt2)^n$.
\end{hint}

\begin{hint}{pb:eggleston}
For (a), restrict $\mu$ to a set where $\mu(D_n(x))\le2^{-nt}$ for all $n\ge N$, and apply the mass distribution principle. For (b), the Bernoulli measure of
parameter $p$ on the binary digits, and the strong law of large numbers. For (c), expand $(a+(1-a))^n$ with $a=\frac kn$.
\end{hint}

\begin{hint}{pb:cantorL}
For (a), take $q=3^{n!}$. For (c), Problem~\ref{pb:wtau}.
\end{hint}

\begin{hint}{pb:weylae}
Expand $|S_N|^2$. Then $\sum_M\int|S_{M^2}|^2<\infty$, and between two squares $S_N$ moves little.
\end{hint}

\begin{hint}{pb:factorial}
For (b), $n!\,e=\sum_{k\le n}\frac{n!}{k!}+\sum_{k>n}\frac{n!}{k!}$. For (c), look at $\sum_k\frac{\varepsilon_k}{k!}$ with $\varepsilon_k\in\{0,1\}$.
\end{hint}

\begin{hint}{pb:del}
For (a), cut $\N$ into blocks $[\rho^m,\rho^{m+1}[$, on which $\sum1/N$ is about $\log\rho$, and choose in each block the $N$ with the smallest $I(N)$. For (b),
compare $S_N$ with $S_{N_{m-1}}$ and let $\rho\to1$ along a sequence.
\end{hint}

\begin{hint}{pb:koksma}
The phase $f(x)=h(x^m-x^n)$ has a monotone derivative of size at least $|h|(m-n)$ on $[a,b]$; integrate by parts. Then apply Problem~\ref{pb:del}.
\end{hint}

\begin{hint}{pb:cantorfourier}
For (a), independence. For (b), $\cos(2\pi3^{m-k})=1$ for $k\le m$. For (c), the map $x\mapsto3x\bmod1$ shifts the ternary digits, so
$\{3^nx\}$ never enters the middle third.
\end{hint}

\begin{hint}{pb:erdosL}
The Liouville numbers form a dense $G_\delta$, and so does $t-\mathcal L$.
\end{hint}

\begin{hint}{pb:halves}
Lebesgue's density theorem.
\end{hint}

\begin{hint}{pb:similar}
For (a), find an interval $I$ in which $A$ has density more than $1-\frac1{2|S|}$, and count the bad $x$ for each point of $S$. For (c), Baire.
\end{hint}

\begin{hint}{pb:typicalcf}
For (a), write the set as a countable intersection of open dense sets, each a union of interiors of cylinders. To make a cylinder bigger in geometric mean, append a
huge digit; to make it smaller, append many ones. For (b), Problem~\ref{pb:khinchin}(b) bounds $a_1\cdots a_n$.
\end{hint}

\begin{hint}{pb:bernsteinset}
For (c), well-order the uncountable closed sets so that each has fewer than $\mathfrak c$ predecessors, and choose two new points in each, by transfinite
recursion. For (d), inner regularity of Lebesgue measure.
\end{hint}

\section*{Solutions}

\begin{solution}{pb:khinconv}
\idea Each fixed $q$ excludes a set of measure about $q\psi(q)$, and the series of these measures converges: that is the whole of the first
Borel--Cantelli lemma. The only care needed is that a solution for $q$ can use only boundedly many numerators.

\writeup (a) Approximations of $x$ and of $x+1$ correspond through $p\mapsto p+q$, so it suffices to treat $x\in[0,1]$, a probability space for
$\lambda$. As $\sum q\psi(q)<\infty$, $q\psi(q)<1$ for $q\ge q_0$. For such $q$, a solution has $|qx-p|<1$, hence $-1\le p\le q+1$; so the set $E_q$ of
$x\in[0,1]$ having a solution with denominator $q$ is covered by $q+3$ intervals of length $2\psi(q)$, and $\lambda(E_q)\le2(q+3)\psi(q)\le8q\psi(q)$.
Since $\sum_q\lambda(E_q)<\infty$, Borel--Cantelli (Proposition~\ref{prop:measureprops}(e)) shows that almost every $x$ lies in finitely many $E_q$. Each $E_q$
with $q\ge q_0$ allows at most $q+3$ numerators, and each $q<q_0$ at most finitely many (as $|p|\le q|x|+q\psi(q)$), so almost every $x$ has finitely many
solutions.

(b) Take $\psi(q)=q^{-2-\varepsilon}$, for which $\sum q^{-1-\varepsilon}<\infty$. A Liouville number $x$ has, for every $n$, a fraction with
$|x-\frac pq|<q^{-n}$ and $q\ge2$; a fixed fraction satisfies this for finitely many $n$ only (as $x$ is irrational), so there are infinitely many fractions
with $|x-\frac pq|<q^{-3}$. By (a) with $\varepsilon=1$ the Liouville numbers form a null set.

(c) With $\psi(q)=\frac1{q^2\log^2q}$ for $q\ge2$, $\sum_q\frac1{q\log^2q}<\infty$. For almost every $x$, which we may also take irrational, only finitely many pairs
$(p_i,q_i)$ satisfy $|x-\frac{p_i}{q_i}|<\psi(q_i)$. If there are no such pairs, take $c(x)=1$. Otherwise put $c(x)=\min\bigl(1,\min_iq_i^2\log^2q_i\,|x-\frac{p_i}{q_i}|\bigr)>0$. Every other pair satisfies
$|x-\frac pq|\ge\frac1{q^2\log^2q}\ge\frac{c(x)}{q^2\log^2q}$, and the exceptional pairs satisfy it by the choice of $c(x)$.

\tools{Chapter~\ref{ch:measures} (Borel--Cantelli); Chapter~\ref{ch:baire} (Liouville numbers).}
\leads Roth's theorem (1955), quoted here for its interest, says that every algebraic irrational $x$ satisfies (b): exponent $2+\varepsilon$ is as good
for $\sqrt2$ as for a random number. Whether $\sqrt[3]2$ satisfies the stronger statement (c), or has bounded digits, nobody knows.
\end{solution}

\begin{solution}{pb:badnull}
\idea ``Badly approximable with constant $c$'' is a closed condition, and it fails at every rational: so it defines a closed set with empty interior.
In the language of continued fractions, $x$ is badly approximable when no digit is large, and a random number has every digit infinitely often.

\writeup (a) For $c>0$ let $\mathrm{Bad}_c$ be the set of real $x$ with $|x-\frac pq|\ge\frac c{q^2}$ for all $p$ and $q\ge1$. It is the complement of a union of
open intervals, hence closed, and it contains no rational number. As $\Q$ is dense, $\mathrm{Bad}_c$ has empty interior: it is nowhere dense. Then
$\mathrm{Bad}=\bigcup_{j\ge1}\mathrm{Bad}_{1/j}$ is meagre.

(b) Suppose $a_n\le K$ for all $n\ge1$. Given $q\ge1$, let $n\ge0$ with $q_n\le q<q_{n+1}$. Then $|qx-p|\ge|q_nx-p_n|=\frac1{x_{n+1}q_n+q_{n-1}}\ge\frac1{(K+2)q_n}\ge\frac1{(K+2)q}$,
since $x_{n+1}<a_{n+1}+1\le K+1$ and $q_{n-1}\le q_n$. So $|x-\frac pq|\ge\frac1{(K+2)q^2}$ and $x\in\mathrm{Bad}$. Conversely, if the digits are unbounded,
$|x-\frac{p_n}{q_n}|<\frac1{a_{n+1}q_n^2}$ with $a_{n+1}$ arbitrarily large, and no $c>0$ works.

(c) $\mathrm{Bad}$ is invariant under $x\mapsto x+1$, so it suffices to look at $(0,1)$. The series $\sum\frac1n$ diverges, so by
Problem~\ref{pb:borelbernstein}(c) with $\varphi(n)=n$, almost every $x\in(0,1)$ has $a_n(x)\ge n$ for infinitely many $n$. (The proof of that problem uses
nothing from this one.) So the digits of almost every $x$ are unbounded, and $\mathrm{Bad}$ is null by (b).

(d) Quadratic irrationals have periodic, hence bounded, digits (Volume~2, Theorem~9.6), so they lie in $\mathrm{Bad}$. More generally, every non-empty open interval of
$(0,1)$ contains a cylinder $I(a_1,\dots,a_n)$ (around an irrational point, cylinders shrink to the point), and that cylinder contains
$[0;a_1,\dots,a_n,1,1,1,\dots]\in\mathrm{Bad}$; so $\mathrm{Bad}$ is dense. The map $(b_n)\mapsto[0;b_1,b_2,\dots]$ is injective on $\{1,2\}^{\N^*}$, which is
uncountable, and its values are badly approximable.

\tools{Volume~2, Chapter~9 (continued fractions); Chapter~\ref{ch:baire} (meagre sets); Problem~\ref{pb:borelbernstein}.}
\leads So $\mathrm{Bad}$ is small for both measure and category. Hausdorff dimension disagrees: Jarník proved in 1928 (quoted) that $\dim_H\mathrm{Bad}=1$, the
dimensions of the sets with digits at most $N$ tending to $1$ (Problem~\ref{pb:etwo} computes one of them). See Bugeaud's book in the sources.
\end{solution}

\begin{solution}{pb:borelbernstein}
\idea The digits of a random number are not independent, but conditioning on the first $n$ digits changes the law of the rest only by a factor between
$\frac12$ and $2$. That is enough for both Borel--Cantelli lemmas: the first needs no independence, and in the second one can multiply the conditional
probabilities of failure step by step.

\writeup (a) A point $x$ of $I$ is $[0;a_1,\dots,a_n,x_{n+1}]$ with $x_{n+1}=\frac1t$ and $t=G^nx\in(0,1)$. By Volume~2, Proposition~9.1(a), this is
$\frac{x_{n+1}p_n+p_{n-1}}{x_{n+1}q_n+q_{n-1}}$, that is $x=\psi(t)=\frac{p_n+tp_{n-1}}{q_n+tq_{n-1}}$. Conversely $\psi(t)$ has first digits $a_1,\dots,a_n$ for every irrational $t\in(0,1)$.
So $I$ is the set of irrational points of the open interval with ends $\psi(0)=\frac{p_n}{q_n}$ and $\psi(1)=\frac{p_n+p_{n-1}}{q_n+q_{n-1}}$, and $\psi$ is a
monotone bijection of $(0,1)$ onto that interval with $|\psi'(t)|=\frac{|p_{n-1}q_n-p_nq_{n-1}|}{(q_n+tq_{n-1})^2}=\frac1{(q_n+tq_{n-1})^2}$. (The numerator is $1$ by Proposition~9.1(b) of the same chapter.) Up to countably many
points, $I\cap G^{-n}B=\psi(B)$, so by the change of variables $x=\psi(t)$ (Theorem~\ref{thm:cov})
\[
\frac{\lambda(I\cap G^{-n}B)}{\lambda(I)}=\frac{\int_B|\psi'|}{\int_0^1|\psi'|},\qquad\frac1{(q_n+q_{n-1})^2}\le|\psi'|\le\frac1{q_n^2},\qquad\int_0^1|\psi'|=\frac1{q_n(q_n+q_{n-1})}.
\]
The ratio is therefore at least $\lambda(B)\frac{q_n}{q_n+q_{n-1}}\ge\frac12\lambda(B)$ and at most $\lambda(B)\frac{q_n+q_{n-1}}{q_n}\le2\lambda(B)$. For irrational
$x\in I$, $a_{n+1}(x)=\lfloor1/G^nx\rfloor\ge k$ if and only if $G^nx<\frac1k$; take $B=(0,\frac1k)$. (For $n=0$, $I=(0,1)$ and the statement is exact.)

(b) The hypothesis forces $\varphi(n)\to\infty$. When $\varphi(n)\ge1$, summing (a) over the cylinders of rank $n-1$ gives
$\lambda(a_n\ge\varphi(n))=\lambda(a_n\ge\lceil\varphi(n)\rceil)\le\frac2{\lceil\varphi(n)\rceil}\le\frac2{\varphi(n)}$, a convergent series; apply Borel--Cantelli.

(c) If $\varphi(n)\le1$ for infinitely many $n$ there is nothing to prove, since $a_n\ge1$. Otherwise $\varphi(n)>1$ for $n\ge n_0$, and $k_n=\lceil\varphi(n)\rceil<2\varphi(n)$,
so $\sum\frac1{k_n}=\infty$. For $n_0\le m<N$ let $C_{m,N}$ be the set where $a_n<k_n$ for $m\le n\le N$. Then $C_{m,N-1}$ is a union of cylinders of rank $N-1$, and on each
of them $a_N\ge k_N$ has conditional probability at least $\frac1{2k_N}$ by (a). Hence $\lambda(C_{m,N})\le\bigl(1-\frac1{2k_N}\bigr)\lambda(C_{m,N-1})$, and by induction
\[
\lambda(C_{m,N})\le\prod_{n=m+1}^N\Bigl(1-\frac1{2k_n}\Bigr)\le\exp\Bigl(-\sum_{n=m+1}^N\frac1{2k_n}\Bigr)\xrightarrow[N\to\infty]{}0.
\]
So $\bigcap_{n\ge m}\{a_n<k_n\}$ is null for every $m$, and almost every $x$ has $a_n\ge k_n\ge\varphi(n)$ for infinitely many $n$.

(d) For every $c>0$, $\sum\frac1{cn\log n}=\infty$, so almost every $x$ has $a_n\ge cn\log n$ infinitely often; intersecting over $c=1,2,3,\dots$ gives
$\limsup\frac{a_n}{n\log n}=\infty$ almost everywhere. And $\sum\frac1{n(\log n)^{1.01}}<\infty$.

\tools{Volume~2, Chapter~9; Chapter~\ref{ch:measures} (Borel--Cantelli); Chapter~\ref{ch:products} (change of variables).}
\leads The estimate in (a) is called \emph{bounded distortion}, and it carries this whole section. Volume~11 uses it again to prove that the Gauss map is ergodic.
What follows is quoted, as a preview of that volume. Large digits make the arithmetic mean of the digits misbehave: it tends to infinity almost everywhere, while the geometric mean converges.
Khinchin showed in 1935 that $\frac{a_1+\dots+a_n}{n\log n}\to\frac1{\log2}$ in measure but not almost everywhere, and Diamond and
Vaaler showed in 1986 that it does converge almost everywhere once the largest digit is removed.
\end{solution}

\begin{solution}{pb:khinchin}
\idea Continued fractions turn ``$x$ is very close to some $\frac pq$'' into ``some digit $a_{n+1}$ is large'', and Problem~\ref{pb:borelbernstein} knows
exactly when digits are large infinitely often. The only obstacle is that $q_n$ depends on $x$. But $q_n$ is at most a product of the digits plus one, the square root of a digit has a finite mean,
and the digits are nearly independent; so $q_n$ grows at most exponentially for almost every $x$. That replaces $q_n$ by the deterministic $e^{5n}$, and the
monotonicity of $q^2\psi(q)$ makes the replacement harmless.

\writeup (a) This is Problem~\ref{pb:khinconv}(a).

(b) From $q_n=a_nq_{n-1}+q_{n-2}\le(a_n+1)q_{n-1}$ and $q_0=1$ we get $q_n\le\prod_{i\le n}(a_i+1)$. Let $P_n(x)=\prod_{i\le n}(a_i(x)+1)^{1/2}$, with $P_0=1$.
On a cylinder $I$ of rank $n-1$ the function $P_{n-1}$ is a constant $P_{n-1}(I)$, and the event $a_n=k$ is $G^{-(n-1)}B$ with $B=\,]\frac1{k+1},\frac1k[$ up to a
countable set, so $\lambda(I\cap\{a_n=k\})\le\frac{2\lambda(I)}{k(k+1)}$ by Problem~\ref{pb:borelbernstein}(a). Hence
\[
\int_IP_n=P_{n-1}(I)\sum_{k\ge1}\sqrt{k+1}\,\lambda(I\cap\{a_n=k\})\le P_{n-1}(I)\lambda(I)\sum_{k\ge1}\frac2{k\sqrt{k+1}}=c\int_IP_{n-1}.
\]
The series converges, since its terms are below $2k^{-3/2}$, and a numerical evaluation gives $c=4.3680\ldots$ The cylinders of rank $n-1$ cover $(0,1)$ up to a countable set, so summing over them gives
$\int_0^1P_n\le c\int_0^1P_{n-1}$, and $\int_0^1P_n\le c^n$ by induction. By Markov's inequality (Proposition~\ref{prop:markov}),
$\lambda(P_n>(2c)^n)\le2^{-n}$, which is summable. By Borel--Cantelli, almost every $x$ has $P_n(x)\le(2c)^n$ for all large $n$, and then
$q_n\le P_n^2\le(2c)^{2n}\le e^{5n}$, because $(2c)^2<77<e^5$.

(c) As $x$ is irrational, $x_{n+1}>a_{n+1}$, so $x_{n+1}q_n+q_{n-1}>a_{n+1}q_n$ and $|x-\frac{p_n}{q_n}|<\frac1{a_{n+1}q_n^2}$. The hypothesis forces $f(q_n)>0$, and then
$\frac1{a_{n+1}q_n^2}\le\frac{f(q_n)}{q_n^2}=\psi(q_n)$.

(d) For $n\ge2$, $Q_{n-1}\ge e^{5n-5}-1$, so $\frac{Q_n}{Q_{n-1}}\le\frac{e^5}{1-e^{-5}}<150$. As $f$ is non-increasing and $\sum_{a<q\le b}\frac1q\le\log\frac ba$,
\[
\sum_{Q_{n-1}<q\le Q_n}q\psi(q)=\sum_{Q_{n-1}<q\le Q_n}\frac{f(q)}q\le f(Q_{n-1})\log150<6f(Q_{n-1}).
\]
Summing over $n\ge2$, $\sum_{q>Q_1}q\psi(q)\le6\sum_{n\ge1}f(Q_n)$, and the divergence of the left side forces that of the right.

(e) Since $\sum f(Q_n)=\infty$ and $f$ is non-increasing, $f>0$ everywhere. Apply Problem~\ref{pb:borelbernstein}(c) with $\varphi(n+1)=1/f(Q_n)$: for almost every
$x$, $a_{n+1}(x)\ge1/f(Q_n)$ for infinitely many $n$. For almost every $x$, (b) also holds; then $q_n\le Q_n$ for large $n$ (an integer at most $e^{5n}$ is at most
$Q_n$), so $f(q_n)\ge f(Q_n)$ and $a_{n+1}\ge1/f(q_n)$ for infinitely many $n$. By (c), $|x-\frac{p_n}{q_n}|<\psi(q_n)$ for these $n$, and distinct $n$ give distinct
fractions. So $\lambda(W(\psi))=1$.

(f) For $\psi(q)=\frac1{q^2\log q}$ when $q\ge3$, with $\psi(1)=\frac1{\log3}$ and $\psi(2)=\frac1{4\log3}$, one has $f(1)=f(2)=f(3)=\frac1{\log3}$ and $f(q)=\frac1{\log q}$ for $q\ge3$, so $f$ is non-increasing. The series $\sum\frac1{q\log q}$ diverges.
For $\psi(q)=\frac1{q^2(\log q)^{1+\varepsilon}}$ the series $\sum\frac1{q(\log q)^{1+\varepsilon}}$ converges.

\tools{Volume~2, Chapter~9; Chapter~\ref{ch:integral} (Markov's inequality); Problems~\ref{pb:khinconv} and \ref{pb:borelbernstein}.}
\leads The constant $5$ in (b) is crude. Lévy proved in 1936 that $\frac1n\log q_n\to\frac{\pi^2}{12\log2}=1.1865\ldots$ for almost every $x$; this is quoted here,
and Volume~11 proves it with the ergodic theorem. The digits stop at $1.1865$, the start of $1.186569\ldots$. Khinchin's theorem does not need the exact constant, only some exponential bound. Without monotonicity the statement is false, as Duffin and Schaeffer showed in 1941. Their conjecture, with reduced fractions and $\varphi(q)\psi(q)$ in place of
$q\psi(q)$, was proved by Koukoulopoulos and Maynard (2020), after Gallagher's zero--one law (1961). Beresnevich and Velani's \emph{mass transference principle}
(2006) turns such Lebesgue statements into Hausdorff ones, and gives Problem~\ref{pb:jarnik} as a special case.
\end{solution}

\begin{solution}{pb:econt}
\idea The quality of the approximations of $e$ is governed by its digits, which grow only linearly, while the denominators grow like a factorial. So
$q^2|e-\frac pq|$ is about $\frac1{a_{n+1}}\approx\frac3{2n}$, which is much larger than $\frac1{\log q}\approx\frac3{n\log n}$: $e$ is approximated \emph{worse}
than a random number.

\writeup (a) From $q_n=a_nq_{n-1}+q_{n-2}$ we get $a_nq_{n-1}\le q_n\le(a_n+1)q_{n-1}$. Hence $q_{3k}\ge\prod_{i\le3k}a_i=\prod_{j=1}^k2j=2^kk!$, and
$q_{3k}\le\prod_{i\le3k}(a_i+1)=2^{2k}\prod_{j=1}^k(2j+1)\le2^{2k}3^kk!\le3^{3k}(k+1)!$. Taking logarithms, $\log q_{3k}=k\log k+O(k)$, and as
$q_{3k}\le q_n\le q_{3k+3}$ for $3k\le n\le3k+3$, $\log q_n=\frac n3\log n+O(n)$.

(b) Let $q\ge3$ and $n$ with $q_n\le q<q_{n+1}$. Then $|qe-p|\ge|q_ne-p_n|>\frac1{(a_{n+1}+2)q_n}\ge\frac1{(a_{n+1}+2)q}$, so
$|e-\frac pq|>\frac1{(a_{n+1}+2)q^2}$. For $m\ge1$ every digit satisfies $a_m\le\frac{2m+2}3$ (the bound fails at $a_0=2$, which this estimate does not use), so $a_{n+1}+2\le n+4$. By (a), $\log q\ge\log q_n\ge c_1n\log n$ for $n\ge2$, hence
$n\le c_2\frac{\log q}{\log\log q}$ for large $q$, and $|e-\frac pq|\ge\frac{c\log\log q}{q^2\log q}$, first for large $q$ and then, decreasing $c$, for all $q\ge3$. As
$\frac{c\log\log q}{\log q}>q^{-\varepsilon}$ for large $q$, the second statement follows.

(c) For large $q$, $\frac{c\log\log q}{q^2\log q}>\frac1{q^2\log q}$, so only finitely many $q$ give solutions, each with finitely many $p$. A numerical calculation lists them. For
$q\ge8$, $\frac1{q^2\log q}<\frac1{2q^2}$, so a solution is a convergent by Legendre's theorem (Volume~2, Theorem~9.5). Testing those convergents, and every integer $p$ for $2\le q<8$, gives nine solutions:
$\frac52$, $\frac62$, $\frac83$, $\frac{11}4$, $\frac{19}7$, $\frac{193}{71}$, $\frac{2721}{1001}$, $\frac{49171}{18089}$, $\frac{1084483}{398959}$.
Eight are in lowest terms. The ninth is $\frac62$, the same value as $3$. Legendre's test does not remove it, because $|e-3|>\frac1{2\cdot2^2}$.

\tools{Volume~2, Chapter~9 (convergents, Legendre's theorem); Problem~\ref{pb:khinchin}.}
\leads Euler found the continued fraction of $e$ in 1737; a short modern proof is due to Cohn (2006). The same method shows that $e$ has irrationality
exponent $2$. For $\pi$, whose digits look random, the best known upper bound for the exponent is about $7.10$ (Zeilberger and Zudilin, 2020), far from the
expected $2$.
\end{solution}

\begin{solution}{pb:wtau}
\idea The natural cover: for each large $q$, the $q+1$ intervals of radius $q^{-\tau}$ around the fractions with denominator $q$. Its $s$-dimensional cost is about
$\sum_qq\cdot q^{-\tau s}$, finite exactly when $s>\frac2\tau$.

\writeup (a) Let $s>\frac2\tau$ and $Q\ge1$. A point of $W(\tau)$ has a solution with $q\ge Q$, and then $-1\le p\le q+1$. So $W(\tau)$ is covered by the intervals of
radius $q^{-\tau}$ around $\frac pq$, $q\ge Q$, $-1\le p\le q+1$, of diameters at most $2Q^{-\tau}$, and
\[
\sum_{q\ge Q}(q+3)(2q^{-\tau})^s\le2^s\cdot4\sum_{q\ge Q}q^{1-\tau s}\xrightarrow[Q\to\infty]{}0,
\]
as $1-\tau s<-1$. So $\mathcal H^s(W(\tau))=0$ and $\dim_HW(\tau)\le s$ (Proposition~\ref{prop:jump}).

(b) As in Problem~\ref{pb:khinconv}(b), a Liouville number in $[0,1]$ lies in every $W(\tau)$, so $\dim_H(\mathcal L\cap[0,1])\le\frac2\tau$ for all $\tau$, hence $0$; and
$\mathcal L$ is a countable union of integer translates of $\mathcal L\cap[0,1]$.

\tools{Chapter~\ref{ch:metricmeasures} (Hausdorff measures).}
\end{solution}

\begin{solution}{pb:jarnik}
\idea The lower bound needs a Cantor set inside $W(\tau)$ that is as fat as possible. Build it greedily: inside each interval of the previous generation, keep a small
interval around every fraction with a much larger denominator. Primes in a dyadic range give fractions that are well spaced, and a very fast growth of the
denominators kills all logarithmic losses.

\writeup (a) Discarding intervals, we may assume each interval of $E_{k-1}$ contains exactly $m_k$ intervals of $E_k$; this shrinks $F$. Since $E_k$ has
$m_1\cdots m_k$ intervals with gaps at least $\varepsilon_k$ inside $[0,1]$, $\varepsilon_k\to0$. Give each interval of $E_k$ the mass $(m_1\cdots m_k)^{-1}$. The discarding step makes a parent contain exactly $m_k$ children, so its mass equals the sum of theirs: the assignment is consistent from one generation to the next. On the algebra of finite unions of these intervals the summed masses are finitely additive, and a countable disjoint union in the algebra is finite because the intervals are closed and separated by gaps. Theorem~\ref{thm:extension} therefore produces a Borel probability $\mu$ carried by $F$ with these masses. (Normalised Lebesgue measure on $E_k$ does not: the intervals have lengths $q^{-\tau}$ for denominators in a dyadic range, so the lengths differ by as much as $2^\tau$ and the masses are not equal.) Let $s$ be less than the $\liminf$, and $k_0$ such that
$m_1\cdots m_{k-1}(m_k\varepsilon_k)^s\ge1$ for $k\ge k_0$. Let $U$ be an interval with $0<|U|<\varepsilon_{k_0}$, and $k>k_0$ with $\varepsilon_k\le|U|<\varepsilon_{k-1}$.
Then $U$ meets at most one interval of $E_{k-1}$, hence at most $m_k$ intervals of $E_k$; and at most $\frac{|U|}{\varepsilon_k}+1\le\frac{2|U|}{\varepsilon_k}$
of them, as they are separated by gaps $\ge\varepsilon_k$. Using $\min(a,b)\le a^{1-s}b^s$,
\[
\mu(U)\le\frac{\min(m_k,2|U|/\varepsilon_k)}{m_1\cdots m_k}\le\frac{m_k^{1-s}(2|U|)^s\varepsilon_k^{-s}}{m_1\cdots m_k}=\frac{2^s|U|^s}{m_1\cdots m_{k-1}(m_k\varepsilon_k)^s}\le2^s|U|^s.
\]
Every closed set of small diameter lies in an interval of the same diameter, so the mass distribution principle (Lemma~\ref{lem:massdist}) gives $\dim_HF\ge s$.

(b) By the prime number theorem (Volume~6, Theorem~5.5) there are at least $\frac Q{2\log Q}$ primes in $]Q,2Q]$ for $Q$ large. For each such prime $q$, $J$ contains at least
$Lq-1\ge LQ-1\ge\frac{LQ}2$ fractions $\frac pq$. As $J\subseteq\,]0,1[$, these fractions are in lowest terms, so different pairs $(p,q)$ give different numbers, and in
total there are at least $\frac{LQ^2}{4\log Q}$. Two of them satisfy $|\frac pq-\frac{p'}{q'}|=\frac{|pq'-p'q|}{qq'}\ge\frac1{4Q^2}$.

(c) The conditions below hold when $Q_k$ is large enough in terms of $Q_{k-1}$, which is the case for $\log Q_k=2^{k^2}$ and $k\ge k_0$; we start the construction at
$k_0$ (renumbering). Let $E_0=[0,1]$ and, by convention, $Q_0^{-\tau}=1$. Given an interval $I$ of $E_{k-1}$, of length $|I|\ge(2Q_{k-1})^{-\tau}$, let $J$ be its middle half, of
length $L=\frac12|I|$, and put into $E_k$ the intervals $[\frac pq-\frac12q^{-\tau},\frac pq+\frac12q^{-\tau}]$ for the fractions $\frac pq\in J$ with $q$ prime in
$]Q_k,2Q_k]$. They lie in $I$ because $\frac12q^{-\tau}\le\frac12Q_k^{-\tau}\le\frac14|I|$. By (b), each interval of $E_{k-1}$ contains
\[
m_k\ge\frac{LQ_k^2}{4\log Q_k}\ge\frac{(2Q_{k-1})^{-\tau}Q_k^2}{8\log Q_k}
\]
of them (we need $LQ_k\ge2$), and any two intervals of $E_k$ are separated by at least $\frac1{4Q_k^2}-Q_k^{-\tau}\ge\frac1{8Q_k^2}=:\varepsilon_k$, as $\tau>2$. A point $x\in F$
lies, for each $k$, in an interval around some $\frac{p_k}{q_k}$ with $q_k\in\,]Q_k,2Q_k]$, so $|x-\frac{p_k}{q_k}|\le\frac12q_k^{-\tau}<q_k^{-\tau}$; the $q_k$ are distinct, the fractions
are in lowest terms, and $x\in W(\tau)$.

(d) From (c), $m_k\varepsilon_k\ge\frac{(2Q_{k-1})^{-\tau}}{64\log Q_k}$, so
\[
\begin{aligned}
-\log(m_k\varepsilon_k)&\le\tau\log Q_{k-1}+\log\log Q_k+O(1),\\
\log(m_1\cdots m_{k-1})&\ge2\log Q_{k-1}-\tau\sum_{j\le k-2}\log Q_j\\
&\quad-\sum_{j\le k-1}\log\log Q_j-O(k).
\end{aligned}
\]
With $\log Q_j=2^{j^2}$: $\sum_{j\le k-2}2^{j^2}\le2\cdot2^{(k-2)^2}$, $\log\log Q_j=j^2\log2$, and both are $o(2^{(k-1)^2})=o(\log Q_{k-1})$. The quotient in (a) tends to
$\frac{2\log Q_{k-1}}{\tau\log Q_{k-1}}=\frac2\tau$. So $\dim_HW(\tau)\ge\dim_HF\ge\frac2\tau$, and Problem~\ref{pb:wtau} gives equality.

\tools{Volume~6, Chapter~5 (prime number theorem); Chapter~\ref{ch:metricmeasures} (mass distribution principle); Problem~\ref{pb:wtau}.}
\leads Jarník proved this in 1929 and Besicovitch independently in 1934. The modern way, quoted here, is the mass transference principle of Beresnevich and Velani (2006), which
deduces the Hausdorff measure of $\limsup$ sets of balls from the Lebesgue measure of $\limsup$ sets of their enlargements: Jarník--Besicovitch then follows from
Dirichlet's theorem.
\end{solution}

\begin{solution}{pb:good}
\idea A cylinder whose digits are all at least $N$ has length at most the product of the inverse squares of its digits, and $\sum_{a\ge N}a^{-2s}$ is small exactly
when $s>\frac12$. Only the first few digits are unconstrained, and they cost a fixed factor.

\writeup (a) $D$ is contained in the union over $m\ge0$ of the sets $D'_m$ where $a_n\ge2$ for all $n>m$. Fix $m$. On every cylinder $I$ of rank $n-1$,
the event $a_n=1$ is $G^{-(n-1)}B$ with $B=\,]\frac12,1[$, so it has conditional probability at least $\frac12\lambda(B)=\frac14$ by Problem~\ref{pb:borelbernstein}(a).
The set where $a_{m+1},\dots,a_{n-1}\ge2$ is a union of such cylinders, so the set where $a_{m+1},\dots,a_n\ge2$ has measure at most $\frac34$ of it, and by induction
at most $(\frac34)^{n-m}$. Letting $n\to\infty$, $\lambda(D'_m)=0$, and $D$ is null. It is uncountable: the numbers
$[0;b_1,b_2,\dots]$ with $b_n\in\{n,n+1\}$ are distinct and lie in $D$.

(b) As $q_n\ge a_nq_{n-1}$, $q_n\ge a_1\cdots a_n$ and $|I(a_1,\dots,a_n)|\le q_n^{-2}\le(a_1\cdots a_n)^{-2}$. Let $s>\frac12$, choose $N$ with
$\theta=\sum_{a\ge N}a^{-2s}<1$, and for $m\ge0$ let $D_m$ be the set of irrational $x$ with $a_n(x)\ge N$ for all $n>m$. For $n>m$, $D_m$ is covered by the cylinders
$I(a_1,\dots,a_n)$ with $a_{m+1},\dots,a_n\ge N$, of diameters at most $2^{-(n-1)}$, and
\[
\sum|I(a_1,\dots,a_n)|^s\le\Bigl(\sum_{a\ge1}a^{-2s}\Bigr)^m\theta^{n-m}=\zeta(2s)^m\theta^{n-m}\xrightarrow[n\to\infty]{}0.
\]
So $\mathcal H^s(D_m)=0$. Every $x\in D$ lies in some $D_m$, so $\mathcal H^s(D)=0$ and $\dim_HD\le s$.

\tools{Volume~2, Chapter~9; Chapter~\ref{ch:metricmeasures}; Problem~\ref{pb:borelbernstein}(a).}
\leads Good proved $\dim_HD=\frac12$ in 1941 (quoted); the lower bound uses a Cantor construction as in Problem~\ref{pb:jarnik}. The value $\frac12$ is the exponent where
$\sum_aa^{-2s}$ stops converging: growing digits make the cylinders behave like the intervals $[\frac1{a+1},\frac1a]$, of length about $a^{-2}$.
\end{solution}

\begin{solution}{pb:etwo}
\idea For the upper bound, cover by the $2^n$ cylinders of rank $n$, whose lengths are at most $q_n^{-2}$ with $q_n$ at least a Fibonacci number. For the lower bound,
feed Problem~\ref{pb:jarnik}(a) the intervals spanned by the pieces of $E_{\{1,2\}}$: each splits in two with a gap, and the gap is a fixed gap transported by a
Möbius map whose derivative is about $q_n^{-2}$.

\writeup Write $E=E_{\{1,2\}}$, and for a word $w=a_1\cdots a_n$ let $\psi_w(t)=\frac{p_n+tp_{n-1}}{q_n+tq_{n-1}}$, so that $E\cap I(w)=\psi_w(E)$.

(a) $E$ is the image of the compact space $\{1,2\}^{\N^*}$ under the map $w\mapsto[0;w]$, which is continuous because words agreeing up to rank $n$ land in one cylinder,
of length at most $2^{-(n-1)}$; so $E$ is compact. Since $[0;a_1,a_2,\dots]$ decreases in $a_1$, increases in $a_2$, and so on, the smallest element of $E$ is
$\alpha=[0;\overline{2,1}]$ and the largest $\beta=[0;\overline{1,2}]$; from $\alpha=\frac1{2+\frac1{1+\alpha}}$ we get $2\alpha^2+2\alpha-1=0$, $\alpha=\frac{\sqrt3-1}2$, and
$\beta=\frac1{1+\alpha}=\sqrt3-1$. $E$ is null because its points have bounded digits, so $E\subseteq\mathrm{Bad}$ (Problem~\ref{pb:badnull}(b) and (c)), and uncountable because the map above is injective.

(b) For digits in $\{1,2\}$, $q_n\ge q_{n-1}+q_{n-2}$, so $q_n\ge F_{n+1}\ge\varphi^{n-1}$. Covering $E$ by the $2^n$ cylinders of rank $n$, for $s>\frac{\log2}{2\log\varphi}$:
$\sum|I(w)|^s\le2^n\varphi^{-2s(n-1)}=\varphi^{2s}\bigl(2\varphi^{-2s}\bigr)^n\to0$.

(c) Let $J_w=\psi_w([\alpha,\beta])$, the interval spanned by $E\cap I(w)$, and $E_k=\bigcup_{|w|=k}J_w$, with $E_0=[0,1]$. The children $J_{w1}$ and $J_{w2}$ of $J_w$ are the images under
$\psi_w$ of $J_1=[\frac1{1+\beta},\frac1{1+\alpha}]=[\frac1{\sqrt3},\sqrt3-1]$ and $J_2=[\frac1{2+\beta},\frac1{2+\alpha}]=[\frac{\sqrt3-1}2,1-\frac1{\sqrt3}]$, which are separated by a gap of
length $g=\frac2{\sqrt3}-1=0.1547\ldots$ Since $|\psi_w'|\ge\frac1{(q_n+q_{n-1})^2}\ge\frac1{4q_n^2}$ and $q_n\le(1+\sqrt2)^n$ (by induction from $q_n\le2q_{n-1}+q_{n-2}$, as
$r=1+\sqrt2$ satisfies $r^2=2r+1$), the two children are separated by at least $\varepsilon_{k}=\frac g4(1+\sqrt2)^{-2(k-1)}$ when $|w|=k-1$; the gaps between intervals of
$E_k$ with different parents are larger. The sets $E_k$ satisfy the hypotheses of Problem~\ref{pb:jarnik}(a) with $m_k=2$, and $\bigcap_kE_k=E$, since the nested intervals
$J_{a_1\cdots a_k}$ shrink to $[0;a_1,a_2,\dots]$. Hence
\[
\dim_HE\ge\lim_{k\to\infty}\frac{(k-1)\log2}{2(k-1)\log(1+\sqrt2)+\log\frac2g}=\frac{\log2}{2\log(1+\sqrt2)}.
\]

(d) The inverse branches of the Gauss map on $E$ are $t\mapsto\frac1{a+t}$, $a\in\{1,2\}$, with derivative of modulus $(a+t)^{-2}$, so by the chain rule
$(\mathcal L_s^n1)(t)=\sum_{|w|=n}|\psi_w'(t)|^s$. As $\frac1{4q_n^2}\le|\psi_w'(t)|\le\frac1{q_n^2}$ and $\frac1{2q_n^2}\le|I(w)|\le\frac1{q_n^2}$, the sum
$\sum_{|w|=n}|I(w)|^s$ is within a factor $4^s$ of $(\mathcal L_s^n1)(t)$, which grows like $\rho(s)^n$, $\rho(s)$ being the spectral radius. If $\rho(s)<1$ the cylinder sums
tend to $0$ and $\dim_HE\le s$, as in (b); if $\rho(s)>1$ they blow up, which suggests $\dim_HE\ge s$. Bowen's formula, quoted here, makes this rigorous: $\dim_HE$ is
the unique $s$ with $\rho(s)=1$.

\tools{Volume~2, Chapter~9; Chapter~\ref{ch:metricmeasures}; Problems~\ref{pb:badnull} and~\ref{pb:jarnik}(a).}
\leads All of this is quoted. Jenkinson and Pollicott (2001) computed $\dim_HE_{\{1,2\}}$ to $25$ digits with this operator; simple collocation gives $16$ digits. Hensley showed that the
set with digits at most $N$ has dimension $1-\frac6{\pi^2N}+O\bigl(\frac{\log N}{N^2}\bigr)$, and such dimensions are the key input in Bourgain and Kontorovich's work
on Zaremba's conjecture (2014).
\end{solution}

\begin{solution}{pb:eggleston}
\idea Dimension counts how many intervals of length $2^{-n}$ a set needs. The numbers with a proportion $p$ of ones need about $\binom n{pn}\approx e^{nH(p)}$ of
them, which gives the upper bound. For the lower bound, spread mass with the coin of bias $p$: it lives on $B_p$, and it gives each relevant dyadic interval
mass about $e^{-nH(p)}$.

\writeup (a) Let $t<s$, and $A_N$ the set of $x\in A$ with $\mu(D_n(x))\le2^{-nt}$ for all $n\ge N$, a Borel set as $D_n(x)$ takes countably many values. The $A_N$ increase
to $A$, so $\mu(A_N)>0$ for some $N$; let $\nu$ be the restriction of $\mu$ to $A_N$. Let $U$ be a closed set of diameter $\delta<2^{-N}$ and $n\ge N$ with
$2^{-n-1}\le\delta<2^{-n}$. Then $U$ meets at most two dyadic intervals $D$ of length $2^{-n}$. If $\nu(D)>0$, $D$ contains a point $y\in A_N$, and $D=D_n(y)$, so
$\nu(D)\le\mu(D_n(y))\le2^{-nt}\le(2\delta)^t$. Hence $\nu(U)\le2^{1+t}\delta^t$, and the mass distribution principle gives $\dim_HA\ge\dim_HA_N\ge t$.

(b) Let $\mu_p$ be the law of $\sum_iX_i2^{-i}$ with $X_i$ independent, $P(X_i=1)=p$ (Theorem~\ref{thm:independent}). It has no atoms, and if $x$ has $k$ ones among its first $n$
binary digits, $\mu_p(D_n(x))=p^k(1-p)^{n-k}$. The strong law of large numbers for these biased coins is proved, by fourth moments, in the solution of
Exercise~\ref{ex:salem}; it says that $\mu_p(B_p)=1$. For $x\in B_p$
\[
-\frac1n\log_2\mu_p(D_n(x))=-\frac kn\log_2p-\Bigl(1-\frac kn\Bigr)\log_2(1-p)\longrightarrow\frac{H(p)}{\log2}.
\]
By (a), $\dim_HB_p\ge\frac{H(p)}{\log2}$.

(c) For $0<k<n$ and $a=\frac kn$: $1=\sum_j\binom nja^j(1-a)^{n-j}\ge\binom nka^k(1-a)^{n-k}=\binom nke^{-nH(a)}$; the cases $k=0,n$ are clear. Let
$t>\frac{H(p)}{\log2}$, $\eta=\frac12(t\log2-H(p))$, and $\delta>0$ with $H(a)\le H(p)+\eta$ when $|a-p|<\delta$. Let $B_{p,N}$ be the set of $x$ whose proportion
$k_n(x)/n$ of ones is within $\delta$ of $p$ for all $n\ge N$; then $B_p\subseteq\bigcup_NB_{p,N}$. For $n\ge N$, $B_{p,N}$ is covered by the closed dyadic intervals of
length $2^{-n}$ with $|\frac kn-p|<\delta$, and
\[
\sum|D|^t\le\sum_{|k-pn|<\delta n}\binom nk2^{-nt}\le(n+1)e^{n(H(p)+\eta)-nt\log2}=(n+1)e^{-n\eta}\longrightarrow0.
\]
So $\mathcal H^t(B_{p,N})=0$ for all $N$, $\mathcal H^t(B_p)=0$, and $\dim_HB_p\le t$.

(d) $\dim_HB_{1/3}=\frac{H(1/3)}{\log2}=\log_23-\frac23=0.9182\ldots$ Lebesgue measure on $[0,1]$ is $\mu_{1/2}$, for which the proportion of ones tends to $\frac12$ almost
surely (Theorem~\ref{thm:borelnormal}), so $\lambda(B_{1/3})=0$. The set needs about $e^{0.6365n}$ intervals of length $2^{-n}$ instead of $2^n=e^{0.6931n}$: exponentially fewer, so measure $0$,
but with an exponent close to $\log2$, so dimension close to $1$.

\tools{Chapter~\ref{ch:measures} (independent digits); Exercise~\ref{ex:salem} (strong law for a biased coin); Chapter~\ref{ch:metricmeasures} (mass distribution principle).}
\leads Besicovitch (1934) and Eggleston (1949) proved this; it is the first example of a \emph{multifractal} decomposition, in which a measure ($\lambda$ here) is
split according to its local behaviour and each piece gets its own dimension. The lemma in (a) is due to Billingsley.
\end{solution}

\begin{solution}{pb:cantorL}
\idea Ternary digits $0$ and $2$ placed only at the positions $k!$ make the truncations of $x$ excellent fractions with denominators $3^{n!}$.

\writeup (a) $x$ has ternary digits $0$ and $2$, so $x\in C$. Let $n\ge1$, $q=3^{n!}\ge3$ and $p=q\sum_{k\le n}2\varepsilon_k3^{-k!}\in\Z$. Since infinitely many $\varepsilon_k$ equal
$1$,
\[
0<x-\frac pq=\sum_{k>n}2\varepsilon_k3^{-k!}\le2\sum_{k>n}3^{-k!}<2\sum_{j\ge(n+1)!}3^{-j}=3\cdot3^{-(n+1)!}=3q^{-n-1}\le q^{-n}.
\]
The middle comparison is strict: the geometric series has exponents that are not factorials. So $x-\frac pq<q^{-n}$, which is what Chapter~\ref{ch:baire} asks of a Liouville number, even at $n=1$, $q=3$, where $3q^{-n-1}=q^{-n}$.
$x$ is irrational: its ternary expansion has infinitely many non-zero digits separated by longer and longer runs of zeros, so it is not eventually periodic.

(b) Distinct sequences give distinct $x$ (ternary expansions with digits $0,2$ are unique), and there are uncountably many sequences with infinitely many ones. $C$ is null
(Exercise~\ref{ex:cantorsum}) and so is $\mathcal L$, yet $C\cap\mathcal L$ is uncountable; and $C\cap\mathcal L\subseteq W(\tau)$ for all $\tau$.

(c) $\dim_H(C\cap\mathcal L)\le\dim_H\mathcal L=0$ by Problem~\ref{pb:wtau}.

\tools{Chapters~\ref{ch:baire} and~\ref{ch:measures}; Problem~\ref{pb:wtau}.}
\leads Whether $C$ contains irrational \emph{algebraic} numbers is a famous open question of Mahler; it is expected that it does not, and that every irrational point of
$C$ is transcendental.
\end{solution}

\begin{solution}{pb:weylae}
\idea On average over $x$, distinct frequencies are orthogonal, so the Weyl sum has mean square exactly $\frac1N$. That is summable along the squares, which gives
almost sure convergence along the squares, and between consecutive squares the sum hardly moves.

\writeup Fix $h\ne0$. Since the $a_n$ are distinct, $\int_0^1e(h(a_m-a_n)x)\dd x=\mathbf 1_{m=n}$ and $\int_0^1|S_N|^2=\frac1{N^2}\sum_{m,n\le N}\mathbf 1_{m=n}=\frac1N$. By
monotone convergence, $\int_0^1\sum_M|S_{M^2}|^2=\sum_M\frac1{M^2}<\infty$, so $\sum_M|S_{M^2}(x)|^2<\infty$ and $S_{M^2}(x)\to0$ for almost every $x$. If
$M^2\le N<(M+1)^2$, then $NS_N-M^2S_{M^2}$ is a sum of $N-M^2\le2M$ terms of modulus $1$, so $|S_N|\le|S_{M^2}|+\frac{2M}N\le|S_{M^2}|+\frac2M$, and $S_N(x)\to0$. The
functions are $1$-periodic in $x$, so this holds for almost every real $x$. Intersecting over the countably many $h\ne0$ and applying Weyl's criterion (Volume~5,
Theorem~6.13) proves the theorem.

\tools{Volume~5, Chapter~6 (Weyl's criterion); Chapter~\ref{ch:integral} (monotone convergence).}
\leads Weyl proved this in 1916, in the memoir where his criterion appears. For $a_n=2^n$ it says that almost every number is normal in base $2$, Borel's theorem of
1909 (Theorem~\ref{thm:borelnormal}), but no algebraic irrational number, $\sqrt2$ for instance, is known to be normal in any base. Baker (1981) showed (quoted) that the discrepancy
is $O\bigl(N^{-1/2}(\log N)^{3/2+\varepsilon}\bigr)$ for almost every $x$, whatever the distinct integers $a_n$.
\end{solution}

\begin{solution}{pb:factorial}
\idea For (b), $e$ is the sum of the reciprocal factorials, so $n!\,e$ is an integer plus a tail of size $\frac1n$. For (c), the same holds for every
$\sum\frac{\varepsilon_k}{k!}$.

\writeup (a) Apply Problem~\ref{pb:weylae} to the distinct integers $a_n=n!$, $n\ge1$.

(b) $n!\,e=\sum_{k\le n}\frac{n!}{k!}+r_n$ with an integer first sum and $r_n=\frac1{n+1}+\frac1{(n+1)(n+2)}+\dots$, so $\frac1{n+1}<r_n\le\sum_{j\ge1}(n+1)^{-j}=\frac1n$. For
$n\ge2$, $\{n!\,e\}=r_n\to0$.

(c) For $\varepsilon=(\varepsilon_k)_{k\ge2}\in\{0,1\}^{\N}$ let $x_\varepsilon=\sum_{k\ge2}\frac{\varepsilon_k}{k!}$. As in (b), $\{n!\,x_\varepsilon\}\le\frac1n$ for $n\ge2$, so
$\{n!\,x_\varepsilon\}\to0$. If $\varepsilon$ and $\varepsilon'$ first differ at $k$, then $|x_\varepsilon-x_{\varepsilon'}|\ge\frac1{k!}-\sum_{j>k}\frac1{j!}>\frac1{k!}-\frac1{k\cdot k!}>0$,
so these numbers are distinct, and uncountably many. If $\{n!\,x\}\to0$ and $r=\frac ab$ is rational, then $n!\,r\in\Z$ for $n\ge b$, so $\{n!(x+r)\}\to0$: the set is
invariant under rational translations, hence dense. It is null by (a).

\tools{Problem~\ref{pb:weylae}.}
\leads The same tail argument proves that $e$ is irrational: if $e=\frac ab$, then $\{n!\,e\}=0$ for $n\ge b$, while $0<r_n<1$ (Fourier's proof). For $\pi$, no such
argument is available, and almost nothing is known about the fractional parts $\{n!\,\pi\}$.
\end{solution}

\begin{solution}{pb:del}
\idea Along a geometric sequence $N_m\approx\rho^m$, one can pick the $N$ with small $I(N)$, because each block $[\rho^m,\rho^{m+1}[$ carries a fixed amount of the
harmonic series. Almost sure convergence along the chosen $N_m$ then costs nothing, but interpolating between them costs up to $1-\rho^{-2}$; letting $\rho\to1$
along a sequence removes that cost.

\writeup (a) Let $B_m=[\rho^m,\rho^{m+1}[\cap\,\N$. For $m\ge m_0$ (with $\rho^m(\rho-1)\ge2$), $\sum_{N\in B_m}\frac1N\ge\frac{\rho^{m+1}-\rho^m-1}{\rho^{m+1}}\ge c_\rho=\frac{1-\rho^{-1}}2$.
Choose $N_m\in B_m$ minimising $I$ on $B_m$; then $c_\rho I(N_m)\le\sum_{N\in B_m}\frac{I(N)}N$, and $\sum_{m\ge m_0}I(N_m)\le c_\rho^{-1}\sum_N\frac{I(N)}N<\infty$.

(b) By monotone convergence, the series $\sum_m|S_{N_m}(x)|^2$ converges, and so $S_{N_m}(x)\to0$, for almost every $x$. For $N$ large, let $m$ be such that $\rho^m\le N<\rho^{m+1}$; then $N_{m-1}<\rho^m\le N$, and
$NS_N-N_{m-1}S_{N_{m-1}}$ is a sum of $N-N_{m-1}$ terms of modulus $1$, so
\[
|S_N|\le|S_{N_{m-1}}|+1-\frac{N_{m-1}}N\le|S_{N_{m-1}}|+1-\frac{\rho^{m-1}}{\rho^{m+1}}.
\]
Hence $\limsup_N|S_N(x)|\le1-\rho^{-2}$ for almost every $x$. Applying this to $\rho=1+\frac1j$ for every $j\ge1$ and intersecting the countably many sets of full measure,
$S_N(x)\to0$ almost everywhere.

(c) The method of squares needs $\sum_MI(M^2)<\infty$, essentially $I(N)=O(N^{-1/2-\delta})$; the criterion only needs, for instance, $I(N)=O\bigl((\log N)^{-1-\delta}\bigr)$, a much slower
decay. In Problem~\ref{pb:weylae}, $I(N)=\frac1N$ and either method works.

\tools{Chapter~\ref{ch:integral} (monotone convergence).}
\leads Davenport, Erdős and LeVeque proved this in 1963, and showed that the condition cannot be weakened much. Their criterion is the standard tool of metric
equidistribution, as in the next problem.
\end{solution}

\begin{solution}{pb:koksma}
\idea For $m\ne n$, the phase $h(x^m-x^n)$ oscillates faster and faster, so the integral of $e(h(x^m-x^n))$ over $[a,b]$ is at most about $\frac1{|h|(m-n)}$.
Summing, $I(N)=O\bigl(\frac{\log N}N\bigr)$, and Problem~\ref{pb:del} applies.

\writeup \emph{A lemma.} Let $f$ be real and $C^2$ on $[a,b]$, with $f'$ monotone and $|f'|\ge\kappa>0$. Then $\bigl|\int_a^be(f)\bigr|\le\frac1{\pi\kappa}$. Indeed $f'$ has
constant sign, $e(f)=\frac1{2\pi if'}\frac{\mathrm d}{\mathrm dx}e(f)$, and integrating by parts,
\[
\Bigl|\int_a^be(f)\Bigr|\le\frac1{2\pi}\Bigl(\frac1{|f'(a)|}+\frac1{|f'(b)|}+\int_a^b\Bigl|\Bigl(\frac1{f'}\Bigr)'\Bigr|\Bigr)=\frac1{2\pi}\cdot2\max\Bigl(\frac1{|f'(a)|},\frac1{|f'(b)|}\Bigr)\le\frac1{\pi\kappa},
\]
since $\frac1{f'}$ is monotone, so that the last integral is $\bigl|\frac1{f'(b)}-\frac1{f'(a)}\bigr|$.

Let $1<a<b$, $h\ne0$, $m>n\ge1$, and $f(x)=h(x^m-x^n)$. Then $f'(x)=hx^{n-1}(mx^{m-n}-n)$, so $|f'|\ge|h|(m-n)$ on $[a,b]$, and
$f''(x)=h\bigl(m(m-1)x^{m-2}-n(n-1)x^{n-2}\bigr)$ has the sign of $h$, so $f'$ is monotone. By the lemma,
\[
\begin{aligned}
I(N)&=\frac1{N^2}\sum_{m,n\le N}\int_a^be\bigl(h(x^m-x^n)\bigr)\dd x\\
&\le\frac{b-a}N+\frac2{N^2}\sum_{1\le n<m\le N}\frac1{\pi|h|(m-n)}\\
&\le\frac{b-a}N+\frac{2(1+\log N)}{\pi|h|N},
\end{aligned}
\]
so $\sum_NI(N)/N<\infty$. By Problem~\ref{pb:del}, $\frac1N\sum_{n\le N}e(hx^n)\to0$ for almost every $x\in[a,b]$. Taking all $h\ne0$ and the intervals $[1+\frac1j,j]$, and applying
Weyl's criterion, $(x^n)$ is equidistributed for almost every $x>1$.

Exceptions: every integer $x\ge2$ (all $x^n$ are integers); $\sqrt2$ (the even powers are integers); the golden ratio $\varphi$, since
$\varphi^n+(-\varphi)^{-n}$ is the $n$-th Lucas number, so $\{\varphi^n\}\to0$ along odd $n$ and $\to1$ along even $n$ (Figure~\ref{fig:koksma}); more generally every
Pisot number.

\tools{Volume~5, Chapter~6 (Weyl's criterion); Problem~\ref{pb:del}.}
\leads Koksma proved this in 1935, for $(x^n)$ and many other families. For no natural number such as $\frac32$ or $e$ is $(x^n)$ known to be equidistributed;
it is not even known whether $\{(3/2)^n\}$ is dense in $[0,1]$. Mahler (1968) called $\xi>0$ a $Z$-number if $\{\xi(3/2)^n\}\in[0,\frac12[$ for all $n\ge0$, showed that each
interval $[k,k+1[$ contains at most one, and asked whether any exist; nobody knows.
\end{solution}

\begin{solution}{pb:cantorfourier}
\idea $X$ is a sum of independent pieces, so its Fourier transform is a product of cosines; multiplying $t$ by $3$ shifts the product and only removes a factor
equal to $1$. For (c), multiplication by $3$ is the shift on the ternary digits, and a shifted point of the Cantor set is still in the Cantor set.

\writeup (a) The partial sums converge to $X$, so by independence and dominated convergence,
\[
\hat\mu(t)=\prod_{k\ge1}\mathbb E\,e(-2t\varepsilon_k3^{-k})=\prod_{k\ge1}\frac{1+e(-2t3^{-k})}2.
\]
With $\frac{1+e^{-i\theta}}2=e^{-i\theta/2}\cos\frac\theta2$ and $\theta=4\pi t3^{-k}$,
each factor is $e(-t3^{-k})\cos(2\pi t3^{-k})$, and $\prod_ke(-t3^{-k})=e(-\frac t2)$.

(b) $|\hat\mu(3^m)|=\prod_{k\ge1}|\cos(2\pi3^{m-k})|$; the factors with $k\le m$ equal $1$, and the others are $\prod_{j\ge1}|\cos(2\pi3^{-j})|=|\hat\mu(1)|$. This product
is not $0$: $\cos\frac{2\pi}3=-\frac12$, and for $j\ge2$ the factors are positive with $\sum_j(1-\cos(2\pi3^{-j}))<\infty$. The Riemann--Lebesgue lemma
(Exercise~\ref{ex:rlemma}) concerns measures with an integrable density; $\mu$ is carried by the null set $C$, so it is singular and the lemma does not apply.

(c) Let $x\in C$, and write $x=\sum_{k\ge1}d_k3^{-k}$ with all $d_k\in\{0,2\}$. Then $3x=d_1+y$ with $y=\sum_{k\ge1}d_{k+1}3^{-k}\in C\subseteq[0,1]$, so
$\{3x\}$ is $y$, or $0$ if $y=1$: in both cases a point of $C$. By induction $\{3^nx\}\in C$ for all $n$. The Cantor set does not meet the open middle third
$]\frac13,\frac23[$, so no term of the sequence falls there, while equidistribution would send a proportion $\frac13$ of the terms there.

\tools{Chapter~\ref{ch:measures} (the Cantor set, independent digits); Chapter~\ref{ch:integral} (dominated convergence, Riemann--Lebesgue).}
\leads Part (c) holds for every point of $C$, and so says nothing of how the orbit is spread inside $C$. The true statement is quoted here as a preview of Volume~11,
where it is a problem on the ergodic theorem: for $\mu$-almost every $x$ the sequence $(3^nx)$ is distributed according to $\mu$ itself. In the other direction,
Cassels (1959) and Schmidt (1960) proved that $\mu$-almost every $x$ is normal in base $2$: the Cantor measure is invisible to $\times3$ but
looks like Lebesgue measure to $\times2$. Hochman and Shmerkin (2015) extended this to all multiplicatively independent bases, a theme that goes back to
Furstenberg's $\times2\times3$ problem.
\end{solution}

\begin{solution}{pb:erdosL}
\idea Liouville numbers are rare for measure but typical for category, and two comeagre sets always meet.

\writeup By Theorem~\ref{thm:liouville}, $\mathcal L$ is a dense $G_\delta$. For $t\in\R$, $y\mapsto t-y$ is a homeomorphism of $\R$, so $t-\mathcal L$ is also a dense $G_\delta$, and
$\mathcal L\cap(t-\mathcal L)$ is comeagre. A countable set is meagre, and $\R$ is not meagre by Baire's theorem (Theorem~\ref{thm:baire}), so this intersection is uncountable.
Each $y$ in it gives $t=y+(t-y)$ with $y$ and $t-y$ Liouville, and each representation comes from at most two such $y$.

\tools{Chapter~\ref{ch:baire} (Baire's theorem, Liouville numbers).}
\leads Erdős published this in 1962, together with the product version: every non-zero real is a product of two Liouville numbers. For sets of positive measure,
$A-A$ contains an interval (Steinhaus, Exercise~\ref{ex:steinhaus}); the Liouville numbers show that a null set of dimension $0$ can do even better.
\end{solution}

\begin{solution}{pb:halves}
\idea Lebesgue's density theorem says that a measurable set fills almost all of every small interval around almost every one of its points.

\writeup Suppose $\varepsilon\lambda(I)\le\lambda(A\cap I)\le(1-\varepsilon)\lambda(I)$ for every interval $I$. Then $\lambda(A\cap[-1,1])\ge2\varepsilon>0$. By Lebesgue's
differentiation theorem (Theorem~\ref{thm:lebdiff}) applied to $\mathbf 1_{A\cap[-1,1]}$, for almost every $x\in A\cap\,]-1,1[$,
$\frac{\lambda(A\cap[x-r,x+r])}{2r}\to1$ as $r\to0$, contradicting the upper bound $1-\varepsilon$. The case $\lambda(A\cap I)=\alpha\lambda(I)$ is the special case
$\varepsilon=\min(\alpha,1-\alpha)$.

\tools{Chapter~\ref{ch:radon} (Lebesgue's differentiation theorem).}
\leads Exercise~\ref{ex:everywhere} builds a set with $0<\lambda(A\cap I)<\lambda(I)$ for every interval: the density of such a set must still oscillate between $0$ and $1$ at
almost every point.
\end{solution}

\begin{solution}{pb:similar}
\idea Near a density point, $A$ fills all but a tiny fraction of an interval, and $|S|$ translates of a tiny set cannot cover half the interval.

\writeup (a) If $|S|\le1$ there is nothing to prove. Otherwise, after an affine change of variables we may assume $S=\{s_1,\dots,s_k\}\subseteq[0,1]$. By the density theorem (Theorem~\ref{thm:lebdiff} applied to $\mathbf 1_A$)
there is an interval $I=[u,u+\ell]$ with $\lambda(A\cap I)>(1-\frac1{2k})\ell$. Let $t=\frac\ell2$. For $x\in[u,u+\frac\ell2]$ all the points $x+ts_i$ lie in $I$. The set
$X_i$ of $x\in[u,u+\frac\ell2]$ with $x+ts_i\notin A$ is a translate of a subset of $I\setminus A$, so $\lambda(X_i)<\frac\ell{2k}$, and $\lambda(\bigcup_iX_i)<\frac\ell2$.
Some $x\in[u,u+\frac\ell2]$ lies in no $X_i$, and then $x+tS\subseteq A$.

(b) $\lambda(A)\ge1-\sum_{n\ge1}2\cdot4^{-n}=\frac13$. $A$ is closed and contains no rational number, so its interior is empty: $A$ is nowhere dense. By (a) it contains an
affine copy of every finite set.

(c) For $t>0$ and $S=\{s_1,s_2,\dots\}$ countable, each $\mathcal L-ts_i$ is a dense $G_\delta$, so $\bigcap_i(\mathcal L-ts_i)$ is comeagre, hence non-empty by Baire's theorem;
any $x$ in it satisfies $x+tS\subseteq\mathcal L$.

\tools{Chapters~\ref{ch:baire} and~\ref{ch:radon}.}
\leads Erdős conjectured in 1974 that (a) fails for every \emph{infinite} $S$: for each infinite $S$ there should be a set of positive measure containing no affine copy
of $S$. Falconer (1984) and Eigen (1985) proved it (quoted) for sequences $s_n\to0$ with $s_{n+1}/s_n\to1$, but the conjecture is open even for $S=\{2^{-n}\}$. By (c), the
analogue for category is false.
\end{solution}

\begin{solution}{pb:typicalcf}
\idea Each condition ``some geometric mean beyond rank $N$ exceeds $M$'' or ``is below $1+\frac1M$'' is open and dense: open because it depends on finitely many digits,
dense because any cylinder can be prolonged by one huge digit or by many ones.

\writeup (a) For $N,M\ge1$ let $U_{N,M}$ be the union of the interiors of the cylinders $I(a_1,\dots,a_n)$ with $n\ge N$ and $(a_1\cdots a_n)^{1/n}>M$, and $V_{N,M}$ the
same with $(a_1\cdots a_n)^{1/n}<1+\frac1M$. These are open. They are dense in $[0,1]$: a non-empty open interval contains an irrational $y$, hence a cylinder
$I(a_1,\dots,a_k)$ around $y$; prolonging the word by one digit $K>M^n$ followed by ones up to length $n=\max(k+1,N)$ gives a sub-cylinder in $U_{N,M}$, and prolonging it
by $n-k$ ones with $n\ge N$ so large that $(a_1\cdots a_k)^{1/n}<1+\frac1M$ gives one in $V_{N,M}$. Let $X'$ be the intersection of all the $U_{N,M}$, $V_{N,M}$ and
$[0,1]\setminus\{r\}$ for $r\in\Q$: a countable intersection of dense open sets, comeagre by Baire's theorem (Theorem~\ref{thm:baire}). A point $x\in X'$ is irrational, and the cylinder of rank $n$ containing
it is $I(a_1(x),\dots,a_n(x))$; so for all $N,M$ there are $n\ge N$ with $(a_1\cdots a_n)^{1/n}>M$ and $n\ge N$ with $(a_1\cdots a_n)^{1/n}<1+\frac1M$. Hence the
$\limsup$ is $\infty$ and the $\liminf$ is at most $1$, and at least $1$ since all digits are $\ge1$. So $X'\subseteq X$, and $X$ is comeagre.

(b) In the solution of Problem~\ref{pb:khinchin}(b), almost every $x$ has $\prod_{i\le n}(a_i+1)\le(2c)^{2n}$ for all large $n$. For such $x$,
$(a_1\cdots a_n)^{1/n}\le(2c)^2<77$ for all large $n$, so the $\limsup$ is finite and $x\notin X$. Hence $X$ is null.

\tools{Chapter~\ref{ch:baire} (Baire's theorem); Problem~\ref{pb:khinchin}(b).}
\leads Much more is true. Khinchin proved in 1935 that for almost every $x$ the geometric mean $(a_1\cdots a_n)^{1/n}$ converges to one and the same constant,
$2.685\ldots$; this is quoted here, and Volume~11 proves it with the ergodic theorem. The argument of (a) shows that the typical number of Baire has no
such constant and no digit frequencies, a pattern already seen for binary digits in
Proposition~\ref{prop:nonnormal}: generic points of Baire are as non-random as possible.
\end{solution}

\begin{solution}{pb:bernsteinset}
\idea There are only $\mathfrak c$ uncountable closed sets, each of size $\mathfrak c$. Listing them so that each comes after fewer than $\mathfrak c$ others, one can
hand each of them a point for $A$ and a point for its complement, never running out of room. A measurable set of positive measure contains an uncountable
closed set, so neither $A$ nor its complement can contain one.

\writeup (a) Let $F'$ be the set of points of $F$ all of whose neighbourhoods meet $F$ in uncountably many points. Each point of $F\setminus F'$ lies in an interval with
rational ends meeting $F$ in a countable set, and there are countably many such intervals, so $F\setminus F'$ is countable and $F'\ne\varnothing$. $F'$ is closed, and
each neighbourhood of a point of $F'$ contains uncountably many points of $F$, hence of $F'$. Choose, for every finite binary word $w$, a closed interval $K_w$ of length at
most $2^{-|w|}$ whose interior meets $F'$, with $K_{w0},K_{w1}$ disjoint inside $K_w$ (possible since the interior of $K_w$ contains two points of $F'$). For every infinite
word, the nested $K_w$ meet $F'$ and shrink to a point of $F'$; distinct words give distinct points. So $|F|\ge|\{0,1\}^\N|=\mathfrak c$ (Theorem~\ref{thm:continuum}), and $|F|\le\mathfrak c$.

(b) A closed set is determined by its complement, which is the union of the intervals with rational ends that it contains; so there are at most $|\mathcal P(\N)|=\mathfrak c$
closed sets, and at least $\mathfrak c$ (the singletons).

(c) The family $\mathcal F$ of uncountable closed sets has cardinality $\mathfrak c$ by (b) (it contains the intervals $[x,x+1]$). By Theorem~\ref{thm:wellordering}, there is a well-order of
$\mathcal F$; if some element has $\mathfrak c$ predecessors, the set $W$ of predecessors of the least such element has cardinality $\mathfrak c$, and each element of $W$ has fewer
than $\mathfrak c$ predecessors. Transporting by a bijection, index $\mathcal F=\{F_\alpha:\alpha\in W\}$ by such a well-ordered set $W$. By transfinite recursion, choose for each
$\alpha$ two distinct points $a_\alpha,b_\alpha\in F_\alpha$ outside $P_\alpha=\{a_\beta,b_\beta:\beta<\alpha\}$; this is possible because $|P_\alpha|<\mathfrak c$ (a set of fewer than
$\mathfrak c$ indices, doubled, still has fewer than $\mathfrak c$ elements) while $|F_\alpha|=\mathfrak c$. Let $A=\{a_\alpha:\alpha\in W\}$. Then $a_\alpha\in F_\alpha\cap A$, and
$b_\alpha\notin A$: $b_\alpha\ne a_\beta$ for $\beta<\alpha$ by the choice of $b_\alpha$, for $\beta=\alpha$ by construction, and for $\beta>\alpha$ by the choice of $a_\beta$.

(d) If $B\subseteq A$ is measurable with $\lambda(B)>0$, it contains a compact $K$ with $\lambda(K)>0$ (Proposition~\ref{prop:regularity}), uncountable since countable sets are null,
so $K$ meets $\R\setminus A$: impossible. The same holds for $\R\setminus A$. If $A$ were measurable, $A\cap[0,1]$ and $[0,1]\setminus A$ would both be null. And if
$\lambda^*(A\cap[0,1])<1$, there would be a Borel $G\supseteq A\cap[0,1]$ with $\lambda(G\cap[0,1])<1$, and $[0,1]\setminus G$ would be a measurable subset of $\R\setminus A$ of
positive measure.

(e) If $A$ were meagre in an interval $J$, then $J\setminus A$ would contain a dense $G_\delta$ set $\bigcap_nU_n$ of $J$, with $U_n$ open and dense in $J$. Choosing as in (a)
closed intervals $K_w$ with $K_w\subseteq U_1\cap\dots\cap U_{|w|}$ produces a Cantor set inside $J\setminus A$, an uncountable closed set missing $A$: impossible. The
complement of $A$ is also a Bernstein set, so $A$ is not comeagre in $J$ either.

\tools{Chapters~\ref{ch:sets} (cardinals, well-ordering), \ref{ch:baire} (Baire) and~\ref{ch:measures} (regularity of Lebesgue measure).}
\leads Felix Bernstein built such sets in 1908. Some choice is unavoidable (quoted): Solovay (1970) constructed a model of set theory without the full axiom of choice, but
with dependent choice, in which every set of reals is Lebesgue measurable, assuming the consistency of an inaccessible cardinal.
\end{solution}

\section*{Map of the chapter}

The last column names the chapters of this volume and the earlier volumes each solution leans on. No solution leans on a later volume.

\begin{center}
\small
\adjustbox{max width=\linewidth}{\begin{tabular}{lll}
Problem & Topic & Leans on\\\hline
\ref{pb:khinconv} & the convergence half of Khinchin's theorem & Ch.~\ref{ch:baire}, \ref{ch:measures}\\
\ref{pb:badnull} & badly approximable numbers: meagre and null & Vol.~2; Ch.~\ref{ch:baire}\\
\ref{pb:borelbernstein} & the Borel--Bernstein theorem on large digits & Vol.~2; Ch.~\ref{ch:measures}, \ref{ch:products}\\
\ref{pb:khinchin} & Khinchin's theorem & Vol.~2; Ch.~\ref{ch:measures}, \ref{ch:integral}\\
\ref{pb:econt} & $e$ is approximated worse than a random number & Vol.~2\\
\ref{pb:wtau} & the dimension of very well approximable numbers: upper bound & Ch.~\ref{ch:metricmeasures}\\
\ref{pb:jarnik} & the Jarník--Besicovitch theorem & Vol.~6; Ch.~\ref{ch:metricmeasures}\\
\ref{pb:good} & numbers whose digits tend to infinity & Vol.~2; Ch.~\ref{ch:metricmeasures}\\
\ref{pb:etwo} & continued fractions with digits $1$ and $2$ & Vol.~2; Ch.~\ref{ch:metricmeasures}\\
\ref{pb:eggleston} & the Besicovitch--Eggleston theorem & Ch.~\ref{ch:measures}, \ref{ch:radon}, \ref{ch:metricmeasures}\\
\ref{pb:cantorL} & Liouville numbers in the Cantor set & Ch.~\ref{ch:baire}, \ref{ch:measures}\\
\ref{pb:weylae} & Weyl's metric theorem & Vol.~5; Ch.~\ref{ch:integral}\\
\ref{pb:factorial} & $n!\,x$ modulo $1$, and $n!\,e$ & Problem~\ref{pb:weylae}\\
\ref{pb:del} & the Davenport--Erdős--LeVeque criterion & Ch.~\ref{ch:integral}\\
\ref{pb:koksma} & Koksma's theorem on $(x^n)$ & Vol.~5\\
\ref{pb:cantorfourier} & the Fourier transform of the Cantor measure & Ch.~\ref{ch:measures}, \ref{ch:integral}\\
\ref{pb:erdosL} & every number is a sum of two Liouville numbers & Ch.~\ref{ch:baire}\\
\ref{pb:halves} & no set halves every interval & Ch.~\ref{ch:radon}\\
\ref{pb:similar} & affine copies in sets of positive measure & Ch.~\ref{ch:baire}, \ref{ch:radon}\\
\ref{pb:typicalcf} & the typical number of Baire against the typical number of Lebesgue & Ch.~\ref{ch:baire}\\
\ref{pb:bernsteinset} & a Bernstein set & Ch.~\ref{ch:sets}, \ref{ch:baire}, \ref{ch:measures}\\
\end{tabular}}
\end{center}

\section*{Sources and further reading}
\addcontentsline{toc}{section}{Sources and further reading}

\begin{itemize}
\item International Mathematical Union, citation for the Fields Medal of James Maynard, 2022 (the epigraph).
\item A.~Khintchine, Einige Sätze über Kettenbrüche, mit Anwendungen auf die Theorie der Diophantischen Approximationen, \emph{Mathematische Annalen} 92 (1924), 115--125;
R.~J.~Duffin and A.~C.~Schaeffer, Khintchine's problem in metric Diophantine approximation, \emph{Duke Mathematical Journal} 8 (1941), 243--255; P.~Gallagher,
Approximation by reduced fractions, \emph{Journal of the Mathematical Society of Japan} 13 (1961), 342--345; D.~Koukoulopoulos and J.~Maynard, On the Duffin--Schaeffer
conjecture, \emph{Annals of Mathematics} 192 (2020), 251--307, arXiv:1907.04593 (the story of the opening).
\item A.~Ya.~Khinchin, \emph{Continued Fractions}, University of Chicago Press, 1964 (Dover reprint 1997); M.~Iosifescu and C.~Kraaikamp, \emph{Metrical Theory of
Continued Fractions}, Kluwer, 2002 (Borel--Bernstein, Khinchin's theorem, Lévy's constant, Khinchin's constant).
\item G.~Harman, \emph{Metric Number Theory}, Oxford University Press, 1998; Y.~Bugeaud, \emph{Approximation by Algebraic Numbers}, Cambridge University Press, 2004
(Jarník's theorem on $\mathrm{Bad}$, Jarník--Besicovitch).
\item V.~Jarník, Diophantische Approximationen und Hausdorffsches Mass, \emph{Mat.\ Sbornik} 36 (1929), 371--382; A.~S.~Besicovitch, Sets of fractional dimensions
(IV): on rational approximation to real numbers, \emph{Journal of the London Mathematical Society} 9 (1934), 126--131; V.~Beresnevich and S.~Velani, A mass transference
principle and the Duffin--Schaeffer conjecture for Hausdorff measures, \emph{Annals of Mathematics} 164 (2006), 971--992.
\item K.~Falconer, \emph{Fractal Geometry: Mathematical Foundations and Applications}, 3rd ed., Wiley, 2014 (the Cantor-type lemma of Problem~\ref{pb:jarnik}, Jarník's theorem,
Besicovitch--Eggleston); I.~J.~Good, The fractional dimensional theory of continued fractions, \emph{Proceedings of the Cambridge Philosophical Society} 37 (1941), 199--228;
H.~G.~Eggleston, The fractional dimension of a set defined by decimal properties, \emph{Quarterly Journal of Mathematics} 20 (1949), 31--36; P.~Billingsley, \emph{Ergodic
Theory and Information}, Wiley, 1965.
\item O.~Jenkinson and M.~Pollicott, Computing the dimension of dynamically defined sets: $E_2$ and bounded continued fractions, \emph{Ergodic Theory and Dynamical Systems} 21
(2001), 1429--1445; D.~Hensley, Continued fraction Cantor sets, Hausdorff dimension, and functional analysis, \emph{Journal of Number Theory} 40 (1992), 336--358;
J.~Bourgain and A.~Kontorovich, On Zaremba's conjecture, \emph{Annals of Mathematics} 180 (2014), 137--196.
\item H.~Cohn, A short proof of the simple continued fraction expansion of $e$, \emph{American Mathematical Monthly} 113 (2006), 57--62; D.~Zeilberger and W.~Zudilin, The
irrationality measure of $\pi$ is at most $7.103205334137\ldots$, \emph{Moscow Journal of Combinatorics and Number Theory} 9 (2020), 407--419.
\item H.~Weyl, Über die Gleichverteilung von Zahlen mod.\ Eins, \emph{Mathematische Annalen} 77 (1916), 313--352; H.~Davenport, P.~Erdős and W.~J.~LeVeque, On Weyl's criterion
for uniform distribution, \emph{Michigan Mathematical Journal} 10 (1963), 311--314; J.~F.~Koksma, Ein mengentheoretischer Satz über die Gleichverteilung modulo Eins,
\emph{Compositio Mathematica} 2 (1935), 250--258; R.~C.~Baker, Metric number theory and the large sieve, \emph{Journal of the London Mathematical Society} 24 (1981), 34--40;
K.~Mahler, An unsolved problem on the powers of $3/2$, \emph{Journal of the Australian Mathematical Society} 8 (1968), 313--321.
\item J.~W.~S.~Cassels, On a problem of Steinhaus about normal numbers, \emph{Colloquium Mathematicum} 7 (1959), 95--101; W.~M.~Schmidt, On normal numbers, \emph{Pacific
Journal of Mathematics} 10 (1960), 661--672; M.~Hochman and P.~Shmerkin, Equidistribution from fractal measures, \emph{Inventiones Mathematicae} 202 (2015), 427--479.
\item P.~Erdős, Representations of real numbers as sums and products of Liouville numbers, \emph{Michigan Mathematical Journal} 9 (1962), 59--60; J.~C.~Oxtoby, \emph{Measure
and Category}, 2nd ed., Springer, 1980 (measure against category, Bernstein sets); K.~J.~Falconer, On a problem of Erdős on sequences and measurable sets, \emph{Proceedings of
the American Mathematical Society} 90 (1984), 77--78; R.~M.~Solovay, A model of set-theory in which every set of reals is Lebesgue measurable, \emph{Annals of Mathematics} 92
(1970), 1--56.
\end{itemize}
